\documentclass[11pt]{article}

\usepackage[margin=1in]{geometry}
\usepackage[T1]{fontenc}
\usepackage[utf8]{inputenc}
\usepackage{lmodern}
\usepackage{microtype}
\usepackage{amsmath,amssymb}
\usepackage{booktabs}
\usepackage{array}
\usepackage{tabularx}
\usepackage{longtable}
\usepackage{graphicx}
\usepackage{caption}
\usepackage{subcaption}
\usepackage{xcolor}
\usepackage{enumitem}
\usepackage{placeins}
\usepackage[numbers,sort&compress]{natbib}
\usepackage[
  colorlinks=true,
  linkcolor=blue!45!black,
  citecolor=blue!45!black,
  urlcolor=blue!55!black
]{hyperref}

\graphicspath{{../public/figures/}{../results/figures/}}

\title{\vspace{-1.0cm}PPO, GRPO, and On-Policy Distillation for Long-Horizon Post-Training}
\author{Limes Labs Research Workstream}
\date{Public draft v0.3, June 18, 2026}

\newcommand{\grpo}{\textsc{GRPO}}
\newcommand{\opd}{\textsc{OPD}}
\newcommand{\opsd}{\textsc{OPSD}}
\newcommand{\ppo}{\textsc{PPO}}
\newcommand{\GRPOm}{\mathrm{GRPO}}
\newcommand{\OPDm}{\mathrm{OPD}}
\newcommand{\PPOm}{\mathrm{PPO}}
\newcommand{\E}{\mathbb{E}}
\newcommand{\KL}{D_{\mathrm{KL}}}
\newcommand{\Ahat}{\widehat{A}}

% Generated by scripts/build_latex_inputs.py; do not edit by hand.
\newcommand{\DeepSeedCount}{20}
\newcommand{\DeepCaseCount}{18}
\newcommand{\DeepMeanCriticWins}{17}
\newcommand{\DeepMeanGroupWins}{1}
\newcommand{\DeepClearCriticCases}{16}
\newcommand{\DeepNearTieCases}{1}
\newcommand{\DeepClearGroupCases}{1}
\newcommand{\DeepMeanDeltaCorrelation}{0.420}



\begin{document}

\maketitle

\begin{abstract}
Critic-free reinforcement learning, especially Group Relative Policy
Optimization (\grpo), is attractive for reasoning-model post-training because
it removes the learned value model and reduces memory pressure. This paper asks
where that trade changes in long-horizon, variable-length, tool-using
trajectories. Our working hypothesis is conditional: critic-based \ppo\ is a
better candidate when token- or state-level temporal credit is observable and
covered, while \grpo\ remains compelling when group reward contrast is reliable
and value modeling is weak or too expensive. \opd\ and \opsd\ are treated as
complementary on-policy distillation objectives, not direct replacements for
policy optimization. We formalize the method taxonomy, cost accounting, and
failure modes, and run a CPU-only toy benchmark with known oracle advantages.
Across a \DeepSeedCount-seed, \DeepCaseCount-case matrix, critic-style
temporal-difference advantages have higher mean oracle-advantage correlation in
\DeepMeanCriticWins\ cases; a confidence-interval reading gives
\DeepClearCriticCases\ clear critic-favorable cases, \DeepNearTieCases\ near
tie, and \DeepClearGroupCases\ clear group-favorable counterexample. The result
is mechanism evidence, not a frontier-model leaderboard claim.
\end{abstract}

\paragraph{Keywords.}
post-training; reinforcement learning from verifiable rewards; PPO; GRPO;
on-policy distillation; credit assignment; long-horizon agents.

\tableofcontents
\clearpage

\section{Introduction}

Long-horizon post-training for language-model agents is not merely long-context
language modeling. A trajectory may include partial plans, tool calls,
retrievals, edits, failed attempts, retries, compacted summaries, and final
verifier feedback. In such traces, the question ``did the rollout succeed?''
is not the same as ``which token or action helped?'' A scalar terminal reward
can be useful for ranking responses, but it can also broadcast the same update
over helpful subgoals, harmful detours, and no-op padding.

This paper is a public Limes Labs workstream on the method question:
\emph{for variable-length, multi-step, tool-heavy trajectories, when does a
critic/value-model setup recover useful temporal credit that a group-relative
setup discards or blurs?} We deliberately avoid centering the work on one base
model. Z.ai's GLM-5.2 report is discussed only as an industry case study: it
describes compacted, variable-length long-horizon rollouts and a move from
group-wise optimization to critic-based \ppo\ with token-level advantages
\citep{zai2026glm52}. That is not independent causal evidence that \ppo\ beats
\grpo; it is a motivating example of the problem structure studied here.

\paragraph{Contributions.}
We make four contributions. First, we separate \ppo, \grpo, \opd, and \opsd\ by
objective, signal granularity, and extra-model cost. Second, we give a cost and
failure-mode checklist for fair post-training comparisons. Third, we introduce
a small deterministic toy environment with oracle advantages and compare
group-relative and critic-style advantage estimators across a multi-seed matrix.
Fourth, we state testable predictions for future neural and agentic
experiments.

\paragraph{Scope discipline.}
The paper makes a narrow claim about credit-assignment mechanisms. It does not
claim that \ppo\ is universally better than \grpo, that \opd\ is weaker than RL,
or that a finite-state toy predicts frontier-model behavior. We use the toy
because the oracle advantage is known, which lets us ask a question that is
usually hidden in language-model training: whether an estimator is aligned with
the token-level causal structure of a trajectory. This distinction matters for
public research. A method can improve a benchmark for reasons unrelated to the
mechanism under discussion, and a method can be mechanistically promising while
still losing under a particular compute budget, implementation, reward model,
or data mixture.

\paragraph{Paper organization.}
Section 2 defines the methods and separates optimization from distillation.
Section 3 reviews related work. Section 4 gives cost accounting. Sections 5--7
define the toy environment, experimental protocol, and results. Sections 8--10
interpret the mechanism evidence, compare method regimes, and state failure
modes. The appendices provide implementation details, case design, full summary
tables, and raw per-seed rows generated directly from the canonical JSON result
artifact.

\section{Definitions and Method Taxonomy}

Let a policy produce a trajectory
\begin{equation}
  \tau = (s_0, a_0, s_1, a_1, \ldots, s_T),
\end{equation}
where a state may include the prompt, partial response, tool outputs, verifier
state, files, memory, or compacted context. Let $R(\tau)$ be an outcome reward,
possibly combined with per-token KL, process rewards, or unit-test feedback.

\subsection{PPO with a critic}

\ppo\ optimizes a clipped policy-gradient surrogate \citep{schulman2017ppo}:
\begin{equation}
L^{\PPOm}(\theta) =
\E_t \left[
  \min\left(
    \rho_t(\theta) A_t,
    \operatorname{clip}(\rho_t(\theta),1-\epsilon,1+\epsilon) A_t
  \right)
\right],
\end{equation}
where
\begin{equation}
  \rho_t(\theta) =
  \frac{\pi_\theta(a_t \mid s_t)}
       {\pi_{\theta_{\mathrm{old}}}(a_t \mid s_t)}.
\end{equation}
In actor-critic use, a value model $V_\phi$ can estimate one-step temporal
difference errors
\begin{equation}
  \delta_t = \mathcal{r}_t + \gamma V_\phi(s_{t+1}) - V_\phi(s_t),
\end{equation}
or generalized advantage estimates \citep{schulman2015gae}:
\begin{equation}
  A_t^{\mathrm{GAE}(\gamma,\lambda)}
  = \sum_{\ell=0}^{\infty}(\gamma\lambda)^\ell \delta_{t+\ell}
  .
\end{equation}
The critic is valuable only if the value target is learnable, sufficiently
covered, and cheaper than the credit it recovers.

\subsection{GRPO}

\grpo\ samples a group of completions for the same prompt and normalizes rewards
inside the group \citep{shao2024deepseekmath,deepseek2025r1}. A simplified
response-level advantage is
\begin{equation}
  \Ahat_i^{\GRPOm} =
  \frac{R_i - \mu_G}{\sigma_G+\epsilon},
\end{equation}
where $\mu_G$ and $\sigma_G$ are group statistics. This critic-free design is
memory-efficient and effective in many verifiable-reward settings. Its
long-horizon risk is intra-rollout information loss: one scalar can be
broadcast across tokens with very different causal roles.

\subsection{OPD and OPSD}

\opd\ trains a student on student-sampled trajectories while a teacher provides
dense supervision on those visited states \citep{agarwal2023gkd,li2026opd}. A
typical objective is
\begin{equation}
L^{\OPDm}
= \E_{x \sim D,\,y \sim \pi_s(\cdot\mid x)}
\left[
  \sum_t
  \KL\!\left(
    \pi_T(\cdot \mid x,y_{<t})
    \,\|\, \pi_s(\cdot \mid x,y_{<t})
  \right)
\right].
\end{equation}
\opsd\ removes the separate teacher model by making one model play a privileged
teacher role and an ordinary student role over student rollouts
\citep{zhao2026opsd}. These are dense distillation objectives. They can be
excellent post-training tools, but they are not the same objective as reward
optimization.

\begin{table}[t]
\centering
\caption{Method taxonomy for long-horizon post-training.}
\label{tab:taxonomy}
\footnotesize
\setlength{\tabcolsep}{3pt}
\begin{tabularx}{\linewidth}{@{}p{2.1cm}p{2.45cm}p{2.05cm}p{2.4cm}X@{}}
\toprule
Method & Objective & Signal granularity & Extra signal & Main risk \\
\midrule
\ppo\ + critic & Reward/KL optimization & token/state & value model, reward, reference & critic cost and miscalibration \\
\grpo & Group-normalized RL & group/response scalar & reward or verifier, no critic & weak intra-rollout credit \\
\opd & Teacher distillation & token logits & teacher model & teacher mismatch and teacher cost \\
\opsd & Privileged self-distillation & token logits & privileged traces & privileged leakage and truncation \\
\bottomrule
\end{tabularx}
\end{table}

\section{Related Work}

\paragraph{PPO and actor--critic post-training.}
\ppo\ was introduced as a simpler clipped-objective alternative to
trust-region policy optimization \citep{schulman2017ppo}. In RLHF-style
language-model training, actor--critic optimization is usually wrapped with KL
penalties, reference policies, reward models, old-logprob storage, reward
normalization, and sequence-level preference signals
\citep{ouyang2022instructgpt,stiennon2020summarize}. The relevant point for
this work is not one historical RLHF recipe, but the availability of a state-
or token-indexed value estimate. A critic can turn delayed outcome feedback
into temporally structured advantages. That can help when rollouts contain
subgoals and detours, but it also creates a second model that can drift,
overfit, or assign confident false credit.

\paragraph{GRPO and RL with verifiable rewards.}
DeepSeekMath introduced \grpo\ as a critic-free method for mathematical
reasoning, replacing the learned value model with group-relative reward
normalization \citep{shao2024deepseekmath}. DeepSeek-R1 then made \grpo-style
RL a central public reference point for reasoning-model training
\citep{deepseek2025r1}. Subsequent work has analyzed the effective loss,
success amplification, response-length effects, dynamic sampling, clipping
details, and statistical structure of \grpo-like gradients
\citep{mroueh2025grpo,liu2025r1zero,yu2025dapo,zhou2026grpo,salmani2026mdpgrpo}.
The practical attraction is real: \grpo\ removes critic memory and critic
training, which can be decisive at large scale. The risk studied here is that a
response-level group scalar can be information-wasteful inside long,
heterogeneous rollouts.

\paragraph{On-policy distillation.}
\opd\ and related generalized knowledge distillation methods train on
student-visited states, reducing train/inference mismatch relative to purely
offline teacher forcing \citep{agarwal2023gkd,li2026opd}. \opsd\ uses the same
model in privileged teacher and ordinary student roles, aiming to capture dense
reasoning supervision without a separate teacher model \citep{zhao2026opsd}.
Recent OPD analyses emphasize teacher/student compatibility, teacher novelty,
length inflation, and truncation collapse \citep{li2026opd,luo2026opd}. These
methods are close to the present problem because they also produce dense
token-level signals on rollouts. They differ because their signal is a
distillation target rather than a reward-improvement objective.

\paragraph{Long-horizon agentic training.}
Long-horizon agentic RL introduces issues that are often absent in short
verifier tasks: tool-call state, file-system state, hidden evaluation
artifacts, environment side effects, compaction, and anti-hacking
infrastructure. The GLM-5.2 report is useful because it explicitly names
compacted sub-traces, variable numbers of trainable traces per prompt, and a
move from group-wise optimization to critic-based \ppo\ with token-level
advantages \citep{zai2026glm52}. We treat this as a case study of a design
pressure, not as proof of a causal performance mechanism.

\section{Cost Accounting}

Fair comparisons should charge all hidden inputs, not only final benchmark
score. For \ppo, the bill includes value-model parameters, value targets,
critic forward passes, optimizer state, old-policy log probabilities, KL
control, and rollout storage. For \grpo, the bill includes group size,
discarded completions, low-dispersion groups, and verifier calls. For
\opd/\opsd, it includes teacher or privileged-context compute, dense logits,
trajectory length inflation, and truncation handling. For agentic coding RL, it
also includes anti-hacking filters, online tool-call monitors, LLM judges,
sandboxing, blocked-call handling, and dummy observations when invalid calls are
intercepted.

The near-term reporting unit for this workstream is:
\begin{equation}
\frac{\Delta \mathrm{reward}}
{\mathrm{generated\ tokens} \times \mathrm{wall\ clock} \times
\mathrm{extra\ model\ memory}},
\end{equation}
reported alongside reproducible artifacts and negative results.

\paragraph{Policy and rollout cost.}
For \ppo, each trainable token typically needs old-policy log probabilities,
advantages, reference-policy information for KL control, and value-model
targets. If rollouts are compacted into sub-traces, the implementation must
decide whether every sub-trace is trainable, how truncation is handled, and how
losses are normalized across very different lengths. \grpo\ avoids the critic
but spends compute on group sampling. If many group members have identical
verifier outcomes, the effective signal can be weak even when token generation
cost is high.

\paragraph{Verifier, judge, and safety cost.}
In coding-agent RL, pass/fail rewards can be corrupted by shortcuts such as
reading hidden tests, fetching target solutions, or exploiting benchmark
leakage. Online anti-hacking filters, rule-based monitors, LLM judges, dummy
observations for blocked calls, and sandboxing are not optional bookkeeping
details. They are part of the training objective as experienced by the agent.
Ignoring them makes method comparisons look cleaner than they really are.

\paragraph{Distillation cost.}
\opd\ and \opsd\ can look cheap if only student updates are counted. A fair bill
also includes teacher logits or hidden states, privileged traces, context
construction, rejection/truncation policies, and any additional sampling needed
to keep the student on-policy. Dense token supervision can be worth this cost,
but it should not be compared to RL methods as if the dense signal were free.

\section{Toy Experiment}

\subsection{Environment}

The toy environment creates variable-length trajectories with three action
types: \texttt{help}, \texttt{harm}, and \texttt{wait}. \texttt{help} increases
a hidden progress score, \texttt{harm} decreases it, and \texttt{wait} consumes
a token without changing progress. A terminal verifier returns success when the
final score reaches a prompt-specific threshold. Because the finite dynamics
are known, we can compute an oracle state-action advantage:
\begin{equation}
  A^\star(s_t,a_t)
  = -c + V^\star(s_{t+1},h-1) - V^\star(s_t,h),
\end{equation}
where $c$ is token cost and $h$ is remaining horizon.

\subsection{Estimators and metrics}

The group-relative estimator normalizes terminal returns inside each prompt
group and broadcasts the resulting scalar to all tokens in the trajectory. The
critic-style estimator fits a tabular value model from sampled returns-to-go and
computes one-step temporal-difference advantages. We measure Pearson
correlation with the oracle advantage, calibrated MSE after affine rescaling,
raw MSE, sign accuracy, wait-token leakage, within-trajectory variance,
zero-variance group fraction, and exact critic-state hit rate.

This is not a closed-loop \ppo-vs-\grpo\ training benchmark. It is an estimator
fidelity experiment under controlled dynamics.

\section{Experimental Protocol}

\paragraph{Design principle.}
The experiment is designed around a known oracle. In real language-model RL, we
usually observe a reward and a model update but not the true token-level
advantage. Here the finite-state dynamics allow exact dynamic programming for
$V^\star$, which turns the question into an estimator-fidelity problem. That
lets us separate two claims that are often conflated: whether a method is cheap
and whether its advantage signal points at the right tokens.

\paragraph{Case axes.}
The deep matrix varies five axes. The horizon axis tests whether scalar
terminal broadcasts become less aligned as trajectories lengthen. The group-size
axis tests whether more completions per prompt rescue \grpo-like normalization.
The critic-budget axis tests coverage: a value table trained from too few
groups should not be trusted. The observability axis removes or coarsens state
features to simulate poor value-model inputs. The sparse-reward axis tests
harder terminal feedback. Finally, the counterexample intentionally combines a
blind critic and low coverage to create a regime where group-relative terminal
outcomes should win.

\paragraph{Decision rule.}
The headline statistic is Pearson correlation with the oracle advantage. A case
is a mean critic win when the mean critic correlation exceeds the mean
group-relative correlation. Because one mean win can be tiny, we also report a
95\% across-seed confidence interval for the delta. A ``clear critic'' case has
positive delta after subtracting the interval; a ``clear group'' case has
negative delta after adding the interval; otherwise the case is labeled a near
tie. This is why the abstract says both ``17 mean critic wins'' and the more
cautious ``16 clear, 1 near tie, 1 clear group'' reading.

\paragraph{Why not a closed-loop toy RL run?}
A closed-loop toy \ppo/\grpo\ implementation is a useful next step, but it
would mix estimator fidelity with optimizer settings, clipping schedules,
policy parameterization, exploration, and KL control. The current paper asks a
more basic question first: when the oracle credit is known, which estimator is
closer to it? That answer should inform later closed-loop experiments rather
than be hidden inside them.

\section{Results}

The canonical matrix uses \DeepSeedCount\ seeds and \DeepCaseCount\ fixed cases
across horizon length, group size, critic budget, observability, sparse-reward
regimes, and a designed counterexample. Critic-style TD has the higher mean
correlation in \DeepMeanCriticWins\ of \DeepCaseCount\ cases. A
confidence-interval reading is more cautious: \DeepClearCriticCases\ cases are
clearly critic-favorable, \DeepNearTieCases\ is a near tie, and
\DeepClearGroupCases\ is clearly group-favorable. The mean critic-minus-group
correlation is \DeepMeanDeltaCorrelation. Table \ref{tab:deep-matrix} reports
representative cases; the full CSV and JSON are committed in \texttt{results/}.

\begin{figure}[t]
\centering
\includegraphics[width=\linewidth]{deep_matrix_delta.png}
\caption{Mean critic-minus-group oracle-advantage correlation by case.
Positive values favor critic-style temporal-difference advantages.}
\label{fig:delta}
\end{figure}

\begin{figure}[t]
\centering
\includegraphics[width=0.86\linewidth]{deep_matrix_coverage.png}
\caption{Critic state coverage versus estimator advantage. The group-relative
counterexample occurs when the critic is blind and undercovered.}
\label{fig:coverage}
\end{figure}

\begin{table}[t]
\centering
\caption{Selected rows from the 20-seed deep matrix. Delta is critic $r$ minus
group $r$; CI is a 95\% across-seed interval for delta.}
\label{tab:deep-matrix}
\small
% Generated by scripts/build_latex_inputs.py; do not edit by hand.
\begin{tabular}{@{}lrrrrl@{}}
\toprule
Case & Group $r$ & Critic $r$ & Delta & CI & Reading \\
\midrule
horizon\_04\_baseline & 0.495 & 0.977 & 0.482 & $\pm$0.011 & critic clear \\
horizon\_16\_long\_wait & 0.293 & 0.865 & 0.571 & $\pm$0.013 & critic clear \\
critic\_budget\_002\_full & 0.352 & 0.367 & 0.015 & $\pm$0.027 & near tie \\
critic\_budget\_128\_full & 0.352 & 0.906 & 0.554 & $\pm$0.017 & critic clear \\
observability\_blind & 0.353 & 0.482 & 0.129 & $\pm$0.017 & critic clear \\
blind\_undercovered\_counterexample & 0.510 & 0.455 & -0.055 & $\pm$0.015 & group clear \\
sparse\_hard\_group08 & 0.310 & 0.910 & 0.600 & $\pm$0.008 & critic clear \\
\bottomrule
\end{tabular}


\end{table}

\paragraph{Interpretation.}
The evidence supports the \ppo-favorable hypothesis in this controlled toy:
when progress state is observable and critic coverage is adequate,
critic-style advantages track token-level oracle credit much better than
response-level group scalars. The result also protects against dogma. In the
coverage-limited case, the confidence interval crosses zero; in the blind and
undercovered case, group-relative terminal outcome information is better than a
poor value table.

\FloatBarrier

\section{Mechanism Analysis}

\paragraph{Horizon length.}
The horizon cases show the central intuition. As trajectories lengthen and wait
tokens accumulate, a response-level scalar is asked to explain increasingly
heterogeneous token roles. The group-relative estimator can still rank complete
rollouts, but that does not mean it assigns good token credit. The critic-style
estimator retains access to progress state and remaining horizon, so it can
separate helpful actions from wait actions even when the final reward is
delayed.

\paragraph{Group size.}
Increasing group size improves the quality of group statistics only when the
additional samples create meaningful reward dispersion. In the long-wait cases,
larger groups help the group-relative estimator somewhat, but the broadcast
problem remains: every token in a completion inherits the same normalized
outcome. In settings where verifier outcomes are sparse or clustered, larger
groups can spend many tokens for little additional advantage information.

\paragraph{Critic budget and coverage.}
The critic-budget axis is the strongest warning against a simplistic ``critics
are better'' claim. With only two training groups, the critic has an exact-state
hit rate below one half and the confidence interval crosses zero. With more
coverage, the critic delta grows sharply. In a neural setting, the analogous
quantity is not table hit rate but value-model calibration under the current
policy distribution. A poorly covered critic should be treated as a risk, not a
badge of sophistication.

\paragraph{Observability.}
The observability cases show that a critic can still extract some value from
coarse state, but a blind critic loses much of the temporal structure. The toy
blind critic is intentionally simple; real long-horizon agents can become
effectively blind for subtler reasons, including lossy compaction, missing tool
state, stale memory, hidden environment variables, or value-model inputs that do
not include the causal feature needed for future return.

\paragraph{Counterexample.}
The counterexample is not an embarrassment; it is the scientific control. If a
paper about \ppo\ versus \grpo\ cannot state when \grpo\ should win, it is not a
method paper. In the blind-undercovered case, terminal group outcomes carry more
usable information than a weak value table. The resulting group win is small but
clear under the confidence-interval rule.

\section{Which Method Is Better?}

The answer is conditional, and the condition is about information. \ppo-style
critic methods are better candidates when future return is predictable from
prefix state, when rollouts contain mixed token roles, and when the critic has
enough coverage to be calibrated. \grpo-style methods are better candidates
when rewards are cheap, verifiable, and meaningfully dispersed within prompt
groups, or when the critic would be blind, stale, or too expensive. \opd\ and
\opsd\ are better candidates when the available dense signal is a teacher or
privileged self-teacher rather than an external reward.

This suggests a hybrid research program. Use rejection sampling and
\grpo-like RL to find successes when verifiers are reliable. Use value or
process models where temporal credit is the bottleneck. Use \opd/\opsd\ to
distill verified or teacher-improved behaviors into cheaper policies. Report
all of these with generated-token, verifier, teacher, critic, and wall-clock
costs. The strongest method is unlikely to be a single objective used
everywhere; it is more likely a pipeline whose components are selected by the
information available at each training stage.

\section{Failure Modes and Testable Predictions}

\paragraph{PPO/critic failures.}
Critics can overfit reward artifacts, become stale after policy drift, collapse
on long sparse trajectories, or assign confident but wrong temporal credit.

\paragraph{GRPO failures.}
\grpo\ can suffer zero-variance groups, low-dispersion amplification, response
length bias, scalar reward broadcast over harmful tokens, and group composition
confounding prompt difficulty with policy quality
\citep{mroueh2025grpo,liu2025r1zero,yu2025dapo,zhou2026grpo,salmani2026mdpgrpo}.

\paragraph{OPD/OPSD failures.}
Teacher/student thinking patterns may be incompatible; the teacher may offer no
new information; on-policy rollouts may inflate in length and truncate; and
privileged self-teacher traces may leak solution structure
\citep{li2026opd,luo2026opd,zhao2026opsd}.

\paragraph{Predictions.}
Increasing horizon length and no-op padding should reduce the alignment between
response-level group advantages and process labels. Increasing group size
should help \grpo\ only when it increases meaningful reward dispersion. Adding
process rewards should narrow the \ppo/\grpo\ credit gap. Degrading critic
observability should produce regimes where \grpo\ wins. \opd/\opsd\ should be
most token-efficient when teacher signals are novel and rollout lengths are
controlled.

\section{Threats to Validity}

\paragraph{Toy dynamics.}
The toy has a small finite state, a known oracle, and a tabular critic. This is
deliberately unlike frontier post-training. The benefit is interpretability;
the cost is that neural generalization, distribution shift, optimizer dynamics,
and reward-model error are not measured.

\paragraph{Estimator versus optimizer.}
The experiment compares advantage estimators, not complete \ppo\ and \grpo\
training loops. A full optimizer could amplify or dampen estimator differences
through clipping, KL schedules, entropy, minibatching, and policy
parameterization. This paper therefore supports a mechanism hypothesis, not a
leaderboard conclusion.

\paragraph{Simplified GRPO baseline.}
The group-relative baseline is intentionally simple. Future comparisons should
include leave-one-out baselines, length-normalized objectives, Dr. GRPO-style
corrections, DAPO-style dynamic sampling, token-level process rewards, and
low-dispersion safeguards. A stronger \grpo\ family member may narrow the gap
in some regimes.

\paragraph{Simplified critic.}
The critic is tabular and trained from returns-to-go. Neural value models can
generalize across states, but they can also hallucinate value structure. The
toy's exact-state hit rate is a proxy for coverage, not a complete model of
critic reliability.

\paragraph{External validity.}
The GLM-5.2 discussion is based on a public vendor report and local text
snapshot, not independent reproduction of its training stack. It should be read
as evidence that the design pressure exists in industry systems, not as proof
that critic-based \ppo\ caused a benchmark gain.

\section{Reproducibility}

The core experiment code has no external Python dependencies. The canonical
checks are:
\begin{verbatim}
python3 -m unittest discover -s tests
./scripts/run_smoke.sh
python3 -m experiments.deep_matrix \
  --output-json results/deep_matrix_20seed.json \
  --output-csv results/deep_matrix_20seed.csv \
  --output-md results/deep_matrix_20seed.md \
  --figures-dir results/figures
./scripts/build_latex_paper.sh
\end{verbatim}
The LaTeX paper is built with Tectonic when available and emitted as the full
paper PDF artifact under \texttt{public/}.

\section{Limitations}

The toy environment has a finite known state and a tabular critic. Real
post-training adds neural approximation, KL controllers, reward-model errors,
tool failures, distributed rollout systems, stale policies, optimizer details,
and data contamination risks. The \grpo\ baseline is intentionally simple; the
next benchmark should include leave-one-out baselines, length-normalized
variants, DAPO/Dr. GRPO-style corrections, process rewards, and dynamic
sampling before making stronger claims.

\section{Roadmap for Limes Labs}

\paragraph{nanoGPT scale.}
The next controlled step is a tiny neural policy where the true environment is
still known but the critic must generalize. This can test whether tabular
coverage results survive function approximation.

\paragraph{AutoResearch and agentic tasks.}
For agentic research workflows, the key artifact is a trace schema that records
tool state, compaction boundaries, blocked actions, reward calls, and process
labels. Without that schema, value modeling and group-relative normalization
cannot be audited fairly.

\paragraph{EuroBench and Parameter Golf.}
Benchmark integrations should report cost-normalized improvement, not only
final score. Parameter Golf is especially useful for comparing memory
multipliers: a critic may improve credit assignment but violate a deployment
budget. EuroBench-style tasks can test whether long-horizon behavior transfers
from toy state to real language and tool use.

\section{Conclusion}

For long heterogeneous trajectories, critic-based \ppo\ is a strong candidate
when value targets are learnable and adequately covered. \grpo\ remains
attractive when group reward contrast is reliable and critic modeling is too
costly or poorly observed. \opd\ and \opsd\ are complementary dense-supervision
tools that can transfer behaviors on student-visited states, but they should
not be conflated with reward optimization. The right public research program is
therefore not ``\ppo\ always beats \grpo,'' but a set of falsifiable boundaries
for when temporal credit is worth its cost.

\clearpage
\appendix

\section{Formal Objective Notes}

\paragraph{Broadcasted group advantages.}
In a response-level group-relative update, a completion-level reward statistic
is often applied to every trainable token in that completion. Let
$g_i=(R_i-\mu_G)/(\sigma_G+\epsilon)$ be the group-normalized scalar for
completion $i$. The token update is then proportional to
$g_i\nabla_\theta\log\pi_\theta(a_{i,t}\mid s_{i,t})$ for many $t$. This is
efficient, but it implicitly assumes that all trainable tokens in the
completion should move in the same reward direction. Long-horizon traces
violate that assumption whenever success depends on a small subset of steps and
the rest are neutral or harmful.

\paragraph{Temporal-difference advantages.}
A critic-style estimator replaces one response-level scalar with a local
difference between predicted future return before and after an action. In the
toy, this is a tabular one-step TD estimate. In language-model agents, the same
idea may be implemented with value heads, separate value models, process reward
models, or learned subgoal predictors. The central requirement is that the
state representation contains enough information to make future return
predictable.

\paragraph{Distillation losses.}
The OPD/OPSD losses in the main text are written as KL objectives because the
teacher signal is a distribution over next tokens. Other implementations may
use sequence-level filtering, hidden-state matching, preference-filtered
teacher traces, or mixed supervised/RL objectives. These variants can be useful
without changing the taxonomy: they are dense imitation or self-imitation
signals rather than direct policy-gradient reward optimization.

\section{Toy Environment Details}

\paragraph{State.}
The state contains prompt identity, hidden progress, remaining horizon, and the
features visible to the critic variant. The full critic observes enough to
distinguish progress states. The coarse critic sees compressed state. The blind
critic intentionally loses the feature needed for reliable future-return
prediction.

\paragraph{Actions.}
The action set is deliberately small. \texttt{help} increases progress,
\texttt{harm} decreases it, and \texttt{wait} consumes budget. This small action
space makes exact dynamic programming feasible and makes the meaning of credit
assignment visible. A language-model trace has richer tokens, but the causal
categories are recognizable: helpful substeps, harmful detours, and neutral
filler.

\paragraph{Reward.}
The terminal verifier is sparse and thresholded. This makes the group-relative
estimator plausible, because rollouts can be ranked by final success, but also
creates the credit problem: the reward is delayed until after many actions.

\paragraph{Oracle advantage.}
The oracle advantage uses exact values under known dynamics. This is the
scientific anchor of the paper. Without it, a method comparison could only say
which estimator produced higher downstream score under one optimizer. With it,
we can ask whether an estimator points in the same direction as the true
token-level credit.

\section{Case Design}

The 18 fixed cases are grouped by axis. Horizon cases increase trace length and
wait-token exposure. Group-size cases vary the number of completions used for
within-prompt normalization. Critic-budget cases vary how much data the value
table sees before evaluation. Observability cases degrade the state available to
the critic. Sparse-reward cases make terminal success harder. The
counterexample combines blindness and low coverage.

\begin{table}[h]
\centering
\caption{Axis-level summary generated from the canonical matrix.}
\small
% Generated by scripts/build_latex_inputs.py; do not edit by hand.
\begin{tabular}{@{}lrrrrrrr@{}}
\toprule
Axis & Cases & Group $r$ & Critic $r$ & Delta & Clear critic & Near tie & Clear group \\
\midrule
counterexample & 1 & 0.510 & 0.455 & -0.055 & 0 & 0 & 1 \\
critic\_budget & 4 & 0.352 & 0.617 & 0.265 & 3 & 1 & 0 \\
group\_size & 4 & 0.305 & 0.868 & 0.563 & 4 & 0 & 0 \\
horizon & 4 & 0.379 & 0.908 & 0.530 & 4 & 0 & 0 \\
observability & 3 & 0.353 & 0.695 & 0.342 & 3 & 0 & 0 \\
sparse\_reward & 2 & 0.296 & 0.878 & 0.582 & 2 & 0 & 0 \\
\bottomrule
\end{tabular}


\end{table}

\section{Full Case Summary}

Table \ref{tab:full-case-summary} gives the full 18-case summary. It is
generated from the canonical matrix JSON rather than hand copied.

\begingroup
\scriptsize
\setlength{\tabcolsep}{3pt}
% Generated by scripts/build_latex_inputs.py; do not edit by hand.
\begin{longtable}{@{}>{\raggedright\arraybackslash}p{3.7cm}>{\raggedright\arraybackslash}p{0.9cm}>{\raggedright\arraybackslash}p{0.8cm}rrrrrr@{}}
\caption{Full 20-seed case summary. Delta is critic $r$ minus group $r$.}\label{tab:full-case-summary}\\
\toprule
Case & Axis & CI read & Group $r$ & Critic $r$ & Delta & CI & Hit & Wait \\
\midrule
\endfirsthead
\toprule
Case & Axis & CI read & Group $r$ & Critic $r$ & Delta & CI & Hit & Wait \\
\midrule
\endhead
horizon 04 baseline & horizon & critic & 0.495 & 0.977 & 0.482 & $\pm$0.011 & 0.999 & 0.303 \\
horizon 08 baseline & horizon & critic & 0.390 & 0.925 & 0.535 & $\pm$0.015 & 0.996 & 0.365 \\
horizon 12 baseline & horizon & critic & 0.335 & 0.866 & 0.531 & $\pm$0.012 & 0.994 & 0.411 \\
horizon 16 long wait & horizon & critic & 0.293 & 0.865 & 0.571 & $\pm$0.013 & 0.996 & 0.539 \\
group 02 long wait & group & critic & 0.213 & 0.771 & 0.557 & $\pm$0.022 & 0.989 & 0.505 \\
group 04 long wait & group & critic & 0.301 & 0.851 & 0.550 & $\pm$0.014 & 0.994 & 0.503 \\
group 08 long wait & group & critic & 0.346 & 0.910 & 0.564 & $\pm$0.011 & 0.997 & 0.503 \\
group 12 long wait & group & critic & 0.358 & 0.940 & 0.582 & $\pm$0.007 & 0.998 & 0.501 \\
budget 002 full & budget & near tie & 0.352 & 0.367 & 0.015 & $\pm$0.027 & 0.476 & 0.388 \\
budget 008 full & budget & critic & 0.352 & 0.459 & 0.107 & $\pm$0.027 & 0.884 & 0.388 \\
budget 032 full & budget & critic & 0.352 & 0.735 & 0.383 & $\pm$0.021 & 0.979 & 0.388 \\
budget 128 full & budget & critic & 0.352 & 0.906 & 0.554 & $\pm$0.017 & 0.996 & 0.388 \\
obs full & obs & critic & 0.352 & 0.864 & 0.513 & $\pm$0.020 & 0.992 & 0.388 \\
obs coarse & obs & critic & 0.356 & 0.738 & 0.382 & $\pm$0.015 & 0.999 & 0.450 \\
obs blind & obs & critic & 0.353 & 0.482 & 0.129 & $\pm$0.017 & 1.000 & 0.444 \\
blind undercovered & counterex. & group & 0.510 & 0.455 & -0.055 & $\pm$0.015 & 0.693 & 0.329 \\
sparse group04 & sparse & critic & 0.283 & 0.846 & 0.563 & $\pm$0.014 & 0.994 & 0.417 \\
sparse group08 & sparse & critic & 0.310 & 0.910 & 0.600 & $\pm$0.008 & 0.997 & 0.420 \\
\bottomrule
\end{longtable}


\endgroup

\section{Raw Seed-Level Rows}

Table \ref{tab:raw-seed-rows} records every case-by-seed row used to compute
the paper's aggregate claims. Publishing these rows makes the mean-win and
confidence-interval interpretation auditable. The table is long by design:
there are \DeepCaseCount\ cases and \DeepSeedCount\ seeds, for 360 rows.

\scriptsize
% Generated by scripts/build_latex_inputs.py; do not edit by hand.
\begin{longtable}{@{}>{\raggedright\arraybackslash}p{4.15cm}rrrrrrr@{}}
\caption{Raw case-by-seed estimator-fidelity rows from the deep matrix.}\label{tab:raw-seed-rows}\\
\toprule
Case & Seed & Group $r$ & Critic $r$ & Delta & Hit & Success & Wait \\
\midrule
\endfirsthead
\toprule
Case & Seed & Group $r$ & Critic $r$ & Delta & Hit & Success & Wait \\
\midrule
\endhead
horizon 04 baseline & 11 & 0.514 & 0.974 & 0.460 & 0.998 & 0.422 & 0.290 \\
horizon 04 baseline & 29 & 0.480 & 0.976 & 0.495 & 0.998 & 0.391 & 0.312 \\
horizon 04 baseline & 47 & 0.477 & 0.982 & 0.505 & 1.000 & 0.396 & 0.286 \\
horizon 04 baseline & 83 & 0.522 & 0.975 & 0.452 & 1.000 & 0.458 & 0.331 \\
horizon 04 baseline & 131 & 0.471 & 0.971 & 0.500 & 1.000 & 0.391 & 0.297 \\
horizon 04 baseline & 173 & 0.505 & 0.971 & 0.466 & 1.000 & 0.422 & 0.296 \\
horizon 04 baseline & 211 & 0.480 & 0.986 & 0.506 & 1.000 & 0.422 & 0.310 \\
horizon 04 baseline & 257 & 0.503 & 0.980 & 0.477 & 1.000 & 0.427 & 0.305 \\
horizon 04 baseline & 307 & 0.493 & 0.985 & 0.492 & 1.000 & 0.438 & 0.317 \\
horizon 04 baseline & 359 & 0.456 & 0.981 & 0.525 & 1.000 & 0.422 & 0.295 \\
horizon 04 baseline & 401 & 0.526 & 0.981 & 0.456 & 1.000 & 0.438 & 0.298 \\
horizon 04 baseline & 463 & 0.491 & 0.981 & 0.490 & 1.000 & 0.417 & 0.309 \\
horizon 04 baseline & 509 & 0.494 & 0.983 & 0.489 & 1.000 & 0.385 & 0.294 \\
horizon 04 baseline & 571 & 0.463 & 0.969 & 0.506 & 0.997 & 0.396 & 0.330 \\
horizon 04 baseline & 631 & 0.471 & 0.977 & 0.506 & 0.997 & 0.474 & 0.265 \\
horizon 04 baseline & 701 & 0.499 & 0.973 & 0.475 & 1.000 & 0.448 & 0.304 \\
horizon 04 baseline & 761 & 0.492 & 0.945 & 0.453 & 0.998 & 0.443 & 0.299 \\
horizon 04 baseline & 823 & 0.548 & 0.974 & 0.427 & 1.000 & 0.464 & 0.312 \\
horizon 04 baseline & 887 & 0.510 & 0.985 & 0.475 & 1.000 & 0.422 & 0.301 \\
horizon 04 baseline & 947 & 0.508 & 0.985 & 0.477 & 1.000 & 0.448 & 0.304 \\
horizon 08 baseline & 11 & 0.410 & 0.898 & 0.488 & 0.995 & 0.542 & 0.361 \\
horizon 08 baseline & 29 & 0.370 & 0.917 & 0.548 & 0.998 & 0.531 & 0.382 \\
horizon 08 baseline & 47 & 0.391 & 0.900 & 0.509 & 0.996 & 0.542 & 0.395 \\
horizon 08 baseline & 83 & 0.365 & 0.923 & 0.558 & 0.996 & 0.490 & 0.369 \\
horizon 08 baseline & 131 & 0.404 & 0.890 & 0.486 & 0.992 & 0.552 & 0.362 \\
horizon 08 baseline & 173 & 0.395 & 0.924 & 0.529 & 0.999 & 0.542 & 0.373 \\
horizon 08 baseline & 211 & 0.402 & 0.919 & 0.517 & 0.994 & 0.547 & 0.347 \\
horizon 08 baseline & 257 & 0.401 & 0.928 & 0.527 & 0.997 & 0.552 & 0.373 \\
horizon 08 baseline & 307 & 0.386 & 0.945 & 0.559 & 1.000 & 0.583 & 0.364 \\
horizon 08 baseline & 359 & 0.339 & 0.957 & 0.618 & 1.000 & 0.526 & 0.376 \\
horizon 08 baseline & 401 & 0.410 & 0.900 & 0.491 & 0.986 & 0.516 & 0.334 \\
horizon 08 baseline & 463 & 0.374 & 0.925 & 0.551 & 0.998 & 0.589 & 0.380 \\
horizon 08 baseline & 509 & 0.390 & 0.919 & 0.529 & 0.998 & 0.599 & 0.369 \\
horizon 08 baseline & 571 & 0.376 & 0.973 & 0.596 & 0.999 & 0.568 & 0.347 \\
horizon 08 baseline & 631 & 0.379 & 0.947 & 0.568 & 1.000 & 0.510 & 0.341 \\
horizon 08 baseline & 701 & 0.400 & 0.925 & 0.526 & 0.996 & 0.583 & 0.404 \\
horizon 08 baseline & 761 & 0.380 & 0.916 & 0.536 & 0.996 & 0.547 & 0.350 \\
horizon 08 baseline & 823 & 0.420 & 0.913 & 0.494 & 0.990 & 0.578 & 0.357 \\
horizon 08 baseline & 887 & 0.408 & 0.937 & 0.530 & 0.995 & 0.604 & 0.364 \\
horizon 08 baseline & 947 & 0.405 & 0.945 & 0.540 & 1.000 & 0.500 & 0.358 \\
horizon 12 baseline & 11 & 0.314 & 0.868 & 0.554 & 0.999 & 0.609 & 0.415 \\
horizon 12 baseline & 29 & 0.348 & 0.919 & 0.571 & 0.995 & 0.589 & 0.400 \\
horizon 12 baseline & 47 & 0.321 & 0.870 & 0.549 & 0.997 & 0.677 & 0.417 \\
horizon 12 baseline & 83 & 0.346 & 0.869 & 0.524 & 0.994 & 0.573 & 0.396 \\
horizon 12 baseline & 131 & 0.339 & 0.878 & 0.540 & 0.994 & 0.604 & 0.397 \\
horizon 12 baseline & 173 & 0.357 & 0.836 & 0.479 & 0.994 & 0.630 & 0.416 \\
horizon 12 baseline & 211 & 0.337 & 0.860 & 0.523 & 0.996 & 0.667 & 0.390 \\
horizon 12 baseline & 257 & 0.346 & 0.893 & 0.547 & 0.997 & 0.656 & 0.437 \\
horizon 12 baseline & 307 & 0.324 & 0.877 & 0.554 & 0.992 & 0.667 & 0.412 \\
horizon 12 baseline & 359 & 0.349 & 0.846 & 0.497 & 0.995 & 0.599 & 0.397 \\
horizon 12 baseline & 401 & 0.349 & 0.870 & 0.521 & 0.993 & 0.630 & 0.412 \\
horizon 12 baseline & 463 & 0.348 & 0.842 & 0.494 & 0.997 & 0.656 & 0.416 \\
horizon 12 baseline & 509 & 0.330 & 0.887 & 0.557 & 0.992 & 0.661 & 0.425 \\
horizon 12 baseline & 571 & 0.326 & 0.840 & 0.514 & 0.986 & 0.661 & 0.424 \\
horizon 12 baseline & 631 & 0.342 & 0.847 & 0.505 & 0.990 & 0.604 & 0.392 \\
horizon 12 baseline & 701 & 0.316 & 0.895 & 0.578 & 0.996 & 0.651 & 0.445 \\
horizon 12 baseline & 761 & 0.349 & 0.884 & 0.536 & 0.990 & 0.625 & 0.405 \\
horizon 12 baseline & 823 & 0.341 & 0.847 & 0.507 & 0.990 & 0.573 & 0.399 \\
horizon 12 baseline & 887 & 0.314 & 0.852 & 0.538 & 0.998 & 0.646 & 0.402 \\
horizon 12 baseline & 947 & 0.314 & 0.844 & 0.530 & 0.993 & 0.620 & 0.421 \\
horizon 16 long wait & 11 & 0.302 & 0.862 & 0.560 & 0.998 & 0.667 & 0.513 \\
horizon 16 long wait & 29 & 0.303 & 0.884 & 0.581 & 0.998 & 0.667 & 0.536 \\
horizon 16 long wait & 47 & 0.266 & 0.866 & 0.601 & 0.998 & 0.718 & 0.538 \\
horizon 16 long wait & 83 & 0.319 & 0.885 & 0.565 & 0.996 & 0.620 & 0.529 \\
horizon 16 long wait & 131 & 0.292 & 0.871 & 0.579 & 0.995 & 0.671 & 0.542 \\
horizon 16 long wait & 173 & 0.297 & 0.869 & 0.572 & 0.997 & 0.606 & 0.552 \\
horizon 16 long wait & 211 & 0.292 & 0.875 & 0.583 & 0.996 & 0.718 & 0.546 \\
horizon 16 long wait & 257 & 0.299 & 0.865 & 0.566 & 0.991 & 0.694 & 0.532 \\
horizon 16 long wait & 307 & 0.264 & 0.909 & 0.645 & 0.998 & 0.671 & 0.539 \\
horizon 16 long wait & 359 & 0.287 & 0.847 & 0.560 & 0.994 & 0.741 & 0.546 \\
horizon 16 long wait & 401 & 0.291 & 0.809 & 0.517 & 0.993 & 0.718 & 0.539 \\
horizon 16 long wait & 463 & 0.299 & 0.860 & 0.561 & 0.997 & 0.662 & 0.552 \\
horizon 16 long wait & 509 & 0.306 & 0.875 & 0.568 & 0.995 & 0.662 & 0.545 \\
horizon 16 long wait & 571 & 0.297 & 0.826 & 0.529 & 0.994 & 0.657 & 0.547 \\
horizon 16 long wait & 631 & 0.305 & 0.840 & 0.534 & 0.995 & 0.667 & 0.517 \\
horizon 16 long wait & 701 & 0.291 & 0.845 & 0.554 & 0.996 & 0.713 & 0.527 \\
horizon 16 long wait & 761 & 0.316 & 0.878 & 0.562 & 0.997 & 0.736 & 0.562 \\
horizon 16 long wait & 823 & 0.287 & 0.866 & 0.579 & 0.995 & 0.657 & 0.532 \\
horizon 16 long wait & 887 & 0.272 & 0.895 & 0.623 & 0.998 & 0.713 & 0.539 \\
horizon 16 long wait & 947 & 0.279 & 0.868 & 0.589 & 0.998 & 0.694 & 0.545 \\
group 02 long wait & 11 & 0.174 & 0.665 & 0.492 & 0.968 & 0.641 & 0.507 \\
group 02 long wait & 29 & 0.227 & 0.774 & 0.547 & 0.991 & 0.719 & 0.496 \\
group 02 long wait & 47 & 0.223 & 0.792 & 0.569 & 0.994 & 0.656 & 0.495 \\
group 02 long wait & 83 & 0.245 & 0.709 & 0.464 & 0.987 & 0.719 & 0.487 \\
group 02 long wait & 131 & 0.189 & 0.758 & 0.569 & 0.998 & 0.625 & 0.514 \\
group 02 long wait & 173 & 0.212 & 0.794 & 0.582 & 0.995 & 0.734 & 0.499 \\
group 02 long wait & 211 & 0.216 & 0.773 & 0.557 & 0.987 & 0.578 & 0.489 \\
group 02 long wait & 257 & 0.236 & 0.775 & 0.539 & 0.986 & 0.641 & 0.515 \\
group 02 long wait & 307 & 0.187 & 0.811 & 0.624 & 0.995 & 0.672 & 0.533 \\
group 02 long wait & 359 & 0.212 & 0.699 & 0.487 & 0.980 & 0.625 & 0.453 \\
group 02 long wait & 401 & 0.195 & 0.720 & 0.525 & 0.996 & 0.672 & 0.528 \\
group 02 long wait & 463 & 0.184 & 0.805 & 0.621 & 0.993 & 0.609 & 0.496 \\
group 02 long wait & 509 & 0.223 & 0.754 & 0.532 & 0.979 & 0.656 & 0.515 \\
group 02 long wait & 571 & 0.247 & 0.878 & 0.631 & 0.986 & 0.656 & 0.536 \\
group 02 long wait & 631 & 0.214 & 0.777 & 0.563 & 0.998 & 0.641 & 0.492 \\
group 02 long wait & 701 & 0.263 & 0.748 & 0.485 & 0.978 & 0.609 & 0.488 \\
group 02 long wait & 761 & 0.225 & 0.783 & 0.558 & 0.998 & 0.672 & 0.490 \\
group 02 long wait & 823 & 0.207 & 0.795 & 0.588 & 0.994 & 0.688 & 0.498 \\
group 02 long wait & 887 & 0.203 & 0.774 & 0.571 & 0.987 & 0.734 & 0.537 \\
group 02 long wait & 947 & 0.187 & 0.827 & 0.641 & 0.991 & 0.703 & 0.542 \\
group 04 long wait & 11 & 0.273 & 0.788 & 0.515 & 0.991 & 0.680 & 0.488 \\
group 04 long wait & 29 & 0.305 & 0.861 & 0.557 & 0.998 & 0.734 & 0.495 \\
group 04 long wait & 47 & 0.265 & 0.794 & 0.529 & 0.988 & 0.656 & 0.481 \\
group 04 long wait & 83 & 0.300 & 0.874 & 0.574 & 1.000 & 0.664 & 0.503 \\
group 04 long wait & 131 & 0.323 & 0.867 & 0.544 & 0.998 & 0.625 & 0.524 \\
group 04 long wait & 173 & 0.333 & 0.796 & 0.464 & 0.985 & 0.656 & 0.465 \\
group 04 long wait & 211 & 0.291 & 0.846 & 0.555 & 0.999 & 0.609 & 0.507 \\
group 04 long wait & 257 & 0.308 & 0.875 & 0.567 & 0.996 & 0.664 & 0.533 \\
group 04 long wait & 307 & 0.331 & 0.879 & 0.548 & 0.991 & 0.641 & 0.495 \\
group 04 long wait & 359 & 0.289 & 0.859 & 0.570 & 0.993 & 0.609 & 0.471 \\
group 04 long wait & 401 & 0.298 & 0.909 & 0.611 & 0.999 & 0.656 & 0.544 \\
group 04 long wait & 463 & 0.293 & 0.880 & 0.587 & 0.995 & 0.570 & 0.487 \\
group 04 long wait & 509 & 0.300 & 0.825 & 0.525 & 0.992 & 0.617 & 0.516 \\
group 04 long wait & 571 & 0.321 & 0.879 & 0.558 & 0.988 & 0.688 & 0.514 \\
group 04 long wait & 631 & 0.303 & 0.841 & 0.538 & 0.998 & 0.680 & 0.485 \\
group 04 long wait & 701 & 0.291 & 0.876 & 0.585 & 0.997 & 0.688 & 0.517 \\
group 04 long wait & 761 & 0.296 & 0.871 & 0.575 & 1.000 & 0.688 & 0.491 \\
group 04 long wait & 823 & 0.296 & 0.857 & 0.561 & 0.982 & 0.664 & 0.508 \\
group 04 long wait & 887 & 0.282 & 0.799 & 0.517 & 0.995 & 0.672 & 0.508 \\
group 04 long wait & 947 & 0.330 & 0.847 & 0.517 & 0.997 & 0.633 & 0.516 \\
group 08 long wait & 11 & 0.337 & 0.903 & 0.566 & 0.997 & 0.656 & 0.506 \\
group 08 long wait & 29 & 0.345 & 0.915 & 0.569 & 0.998 & 0.652 & 0.493 \\
group 08 long wait & 47 & 0.301 & 0.907 & 0.607 & 0.998 & 0.680 & 0.506 \\
group 08 long wait & 83 & 0.356 & 0.931 & 0.575 & 0.997 & 0.625 & 0.460 \\
group 08 long wait & 131 & 0.373 & 0.931 & 0.558 & 0.996 & 0.637 & 0.522 \\
group 08 long wait & 173 & 0.353 & 0.865 & 0.512 & 0.996 & 0.633 & 0.497 \\
group 08 long wait & 211 & 0.357 & 0.927 & 0.570 & 0.998 & 0.621 & 0.503 \\
group 08 long wait & 257 & 0.344 & 0.938 & 0.595 & 0.998 & 0.617 & 0.505 \\
group 08 long wait & 307 & 0.333 & 0.925 & 0.592 & 0.998 & 0.664 & 0.511 \\
group 08 long wait & 359 & 0.350 & 0.891 & 0.541 & 0.995 & 0.660 & 0.505 \\
group 08 long wait & 401 & 0.343 & 0.924 & 0.581 & 0.998 & 0.660 & 0.515 \\
group 08 long wait & 463 & 0.340 & 0.900 & 0.560 & 0.999 & 0.609 & 0.518 \\
group 08 long wait & 509 & 0.333 & 0.925 & 0.592 & 0.996 & 0.691 & 0.503 \\
group 08 long wait & 571 & 0.333 & 0.867 & 0.534 & 0.994 & 0.672 & 0.526 \\
group 08 long wait & 631 & 0.366 & 0.904 & 0.538 & 0.997 & 0.609 & 0.517 \\
group 08 long wait & 701 & 0.342 & 0.899 & 0.556 & 0.998 & 0.668 & 0.505 \\
group 08 long wait & 761 & 0.343 & 0.945 & 0.602 & 0.998 & 0.586 & 0.494 \\
group 08 long wait & 823 & 0.346 & 0.887 & 0.541 & 0.989 & 0.617 & 0.484 \\
group 08 long wait & 887 & 0.369 & 0.914 & 0.545 & 0.997 & 0.613 & 0.484 \\
group 08 long wait & 947 & 0.350 & 0.900 & 0.550 & 0.996 & 0.660 & 0.513 \\
group 12 long wait & 11 & 0.365 & 0.924 & 0.559 & 0.999 & 0.646 & 0.496 \\
group 12 long wait & 29 & 0.367 & 0.935 & 0.568 & 0.999 & 0.633 & 0.491 \\
group 12 long wait & 47 & 0.338 & 0.958 & 0.620 & 0.999 & 0.659 & 0.511 \\
group 12 long wait & 83 & 0.372 & 0.939 & 0.568 & 0.999 & 0.674 & 0.498 \\
group 12 long wait & 131 & 0.353 & 0.958 & 0.605 & 1.000 & 0.638 & 0.531 \\
group 12 long wait & 173 & 0.355 & 0.948 & 0.593 & 0.998 & 0.690 & 0.490 \\
group 12 long wait & 211 & 0.351 & 0.948 & 0.596 & 0.998 & 0.617 & 0.488 \\
group 12 long wait & 257 & 0.352 & 0.940 & 0.588 & 0.999 & 0.607 & 0.511 \\
group 12 long wait & 307 & 0.367 & 0.943 & 0.576 & 1.000 & 0.654 & 0.502 \\
group 12 long wait & 359 & 0.348 & 0.938 & 0.590 & 0.998 & 0.667 & 0.493 \\
group 12 long wait & 401 & 0.363 & 0.937 & 0.574 & 0.997 & 0.674 & 0.502 \\
group 12 long wait & 463 & 0.356 & 0.919 & 0.564 & 0.997 & 0.646 & 0.520 \\
group 12 long wait & 509 & 0.352 & 0.942 & 0.590 & 0.996 & 0.635 & 0.490 \\
group 12 long wait & 571 & 0.368 & 0.928 & 0.560 & 0.996 & 0.628 & 0.511 \\
group 12 long wait & 631 & 0.371 & 0.937 & 0.566 & 0.998 & 0.656 & 0.507 \\
group 12 long wait & 701 & 0.367 & 0.945 & 0.578 & 1.000 & 0.693 & 0.513 \\
group 12 long wait & 761 & 0.344 & 0.934 & 0.590 & 0.998 & 0.659 & 0.500 \\
group 12 long wait & 823 & 0.348 & 0.947 & 0.599 & 0.997 & 0.617 & 0.486 \\
group 12 long wait & 887 & 0.356 & 0.942 & 0.585 & 0.998 & 0.635 & 0.496 \\
group 12 long wait & 947 & 0.363 & 0.942 & 0.579 & 0.999 & 0.656 & 0.494 \\
budget 002 full & 11 & 0.358 & 0.415 & 0.057 & 0.467 & 0.573 & 0.392 \\
budget 002 full & 29 & 0.395 & 0.405 & 0.009 & 0.476 & 0.620 & 0.400 \\
budget 002 full & 47 & 0.342 & 0.330 & -0.012 & 0.536 & 0.609 & 0.385 \\
budget 002 full & 83 & 0.308 & 0.251 & -0.056 & 0.439 & 0.630 & 0.409 \\
budget 002 full & 131 & 0.319 & 0.340 & 0.021 & 0.369 & 0.646 & 0.381 \\
budget 002 full & 173 & 0.328 & 0.441 & 0.113 & 0.514 & 0.562 & 0.381 \\
budget 002 full & 211 & 0.363 & 0.299 & -0.064 & 0.501 & 0.630 & 0.398 \\
budget 002 full & 257 & 0.359 & 0.451 & 0.091 & 0.459 & 0.516 & 0.375 \\
budget 002 full & 307 & 0.357 & 0.267 & -0.090 & 0.522 & 0.594 & 0.380 \\
budget 002 full & 359 & 0.313 & 0.414 & 0.101 & 0.376 & 0.568 & 0.364 \\
budget 002 full & 401 & 0.385 & 0.304 & -0.081 & 0.506 & 0.573 & 0.374 \\
budget 002 full & 463 & 0.381 & 0.484 & 0.102 & 0.454 & 0.557 & 0.407 \\
budget 002 full & 509 & 0.329 & 0.358 & 0.029 & 0.489 & 0.604 & 0.381 \\
budget 002 full & 571 & 0.381 & 0.379 & -0.003 & 0.432 & 0.599 & 0.397 \\
budget 002 full & 631 & 0.364 & 0.306 & -0.058 & 0.470 & 0.562 & 0.388 \\
budget 002 full & 701 & 0.318 & 0.368 & 0.050 & 0.543 & 0.594 & 0.406 \\
budget 002 full & 761 & 0.329 & 0.375 & 0.046 & 0.503 & 0.651 & 0.413 \\
budget 002 full & 823 & 0.387 & 0.409 & 0.022 & 0.389 & 0.615 & 0.392 \\
budget 002 full & 887 & 0.377 & 0.368 & -0.009 & 0.528 & 0.609 & 0.360 \\
budget 002 full & 947 & 0.342 & 0.383 & 0.040 & 0.539 & 0.599 & 0.372 \\
budget 008 full & 11 & 0.358 & 0.471 & 0.114 & 0.868 & 0.573 & 0.392 \\
budget 008 full & 29 & 0.395 & 0.536 & 0.141 & 0.873 & 0.620 & 0.400 \\
budget 008 full & 47 & 0.342 & 0.507 & 0.165 & 0.913 & 0.609 & 0.385 \\
budget 008 full & 83 & 0.308 & 0.470 & 0.163 & 0.856 & 0.630 & 0.409 \\
budget 008 full & 131 & 0.319 & 0.395 & 0.076 & 0.878 & 0.646 & 0.381 \\
budget 008 full & 173 & 0.328 & 0.567 & 0.239 & 0.885 & 0.562 & 0.381 \\
budget 008 full & 211 & 0.363 & 0.466 & 0.104 & 0.846 & 0.630 & 0.398 \\
budget 008 full & 257 & 0.359 & 0.517 & 0.158 & 0.917 & 0.516 & 0.375 \\
budget 008 full & 307 & 0.357 & 0.399 & 0.042 & 0.926 & 0.594 & 0.380 \\
budget 008 full & 359 & 0.313 & 0.530 & 0.217 & 0.872 & 0.568 & 0.364 \\
budget 008 full & 401 & 0.385 & 0.480 & 0.096 & 0.885 & 0.573 & 0.374 \\
budget 008 full & 463 & 0.381 & 0.455 & 0.074 & 0.958 & 0.557 & 0.407 \\
budget 008 full & 509 & 0.329 & 0.435 & 0.106 & 0.913 & 0.604 & 0.381 \\
budget 008 full & 571 & 0.381 & 0.465 & 0.084 & 0.875 & 0.599 & 0.397 \\
budget 008 full & 631 & 0.364 & 0.353 & -0.012 & 0.828 & 0.562 & 0.388 \\
budget 008 full & 701 & 0.318 & 0.422 & 0.104 & 0.884 & 0.594 & 0.406 \\
budget 008 full & 761 & 0.329 & 0.348 & 0.020 & 0.847 & 0.651 & 0.413 \\
budget 008 full & 823 & 0.387 & 0.429 & 0.042 & 0.849 & 0.615 & 0.392 \\
budget 008 full & 887 & 0.377 & 0.501 & 0.124 & 0.903 & 0.609 & 0.360 \\
budget 008 full & 947 & 0.342 & 0.434 & 0.092 & 0.902 & 0.599 & 0.372 \\
budget 032 full & 11 & 0.358 & 0.715 & 0.357 & 0.975 & 0.573 & 0.392 \\
budget 032 full & 29 & 0.395 & 0.761 & 0.365 & 0.974 & 0.620 & 0.400 \\
budget 032 full & 47 & 0.342 & 0.749 & 0.407 & 0.987 & 0.609 & 0.385 \\
budget 032 full & 83 & 0.308 & 0.709 & 0.402 & 0.975 & 0.630 & 0.409 \\
budget 032 full & 131 & 0.319 & 0.715 & 0.396 & 0.974 & 0.646 & 0.381 \\
budget 032 full & 173 & 0.328 & 0.828 & 0.500 & 0.980 & 0.562 & 0.381 \\
budget 032 full & 211 & 0.363 & 0.667 & 0.305 & 0.971 & 0.630 & 0.398 \\
budget 032 full & 257 & 0.359 & 0.795 & 0.436 & 0.993 & 0.516 & 0.375 \\
budget 032 full & 307 & 0.357 & 0.708 & 0.351 & 0.983 & 0.594 & 0.380 \\
budget 032 full & 359 & 0.313 & 0.780 & 0.467 & 0.982 & 0.568 & 0.364 \\
budget 032 full & 401 & 0.385 & 0.691 & 0.306 & 0.973 & 0.573 & 0.374 \\
budget 032 full & 463 & 0.381 & 0.757 & 0.375 & 0.987 & 0.557 & 0.407 \\
budget 032 full & 509 & 0.329 & 0.672 & 0.343 & 0.981 & 0.604 & 0.381 \\
budget 032 full & 571 & 0.381 & 0.717 & 0.336 & 0.979 & 0.599 & 0.397 \\
budget 032 full & 631 & 0.364 & 0.720 & 0.356 & 0.961 & 0.562 & 0.388 \\
budget 032 full & 701 & 0.318 & 0.700 & 0.382 & 0.969 & 0.594 & 0.406 \\
budget 032 full & 761 & 0.329 & 0.700 & 0.372 & 0.973 & 0.651 & 0.413 \\
budget 032 full & 823 & 0.387 & 0.798 & 0.411 & 0.993 & 0.615 & 0.392 \\
budget 032 full & 887 & 0.377 & 0.771 & 0.394 & 0.985 & 0.609 & 0.360 \\
budget 032 full & 947 & 0.342 & 0.740 & 0.397 & 0.984 & 0.599 & 0.372 \\
budget 128 full & 11 & 0.358 & 0.907 & 0.549 & 0.994 & 0.573 & 0.392 \\
budget 128 full & 29 & 0.395 & 0.884 & 0.489 & 0.996 & 0.620 & 0.400 \\
budget 128 full & 47 & 0.342 & 0.909 & 0.567 & 0.997 & 0.609 & 0.385 \\
budget 128 full & 83 & 0.308 & 0.896 & 0.588 & 0.996 & 0.630 & 0.409 \\
budget 128 full & 131 & 0.319 & 0.899 & 0.580 & 0.992 & 0.646 & 0.381 \\
budget 128 full & 173 & 0.328 & 0.942 & 0.614 & 1.000 & 0.562 & 0.381 \\
budget 128 full & 211 & 0.363 & 0.873 & 0.510 & 0.993 & 0.630 & 0.398 \\
budget 128 full & 257 & 0.359 & 0.917 & 0.557 & 0.998 & 0.516 & 0.375 \\
budget 128 full & 307 & 0.357 & 0.880 & 0.523 & 0.993 & 0.594 & 0.380 \\
budget 128 full & 359 & 0.313 & 0.938 & 0.625 & 0.997 & 0.568 & 0.364 \\
budget 128 full & 401 & 0.385 & 0.905 & 0.521 & 0.992 & 0.573 & 0.374 \\
budget 128 full & 463 & 0.381 & 0.915 & 0.533 & 0.996 & 0.557 & 0.407 \\
budget 128 full & 509 & 0.329 & 0.923 & 0.594 & 0.993 & 0.604 & 0.381 \\
budget 128 full & 571 & 0.381 & 0.880 & 0.498 & 0.997 & 0.599 & 0.397 \\
budget 128 full & 631 & 0.364 & 0.881 & 0.516 & 0.987 & 0.562 & 0.388 \\
budget 128 full & 701 & 0.318 & 0.909 & 0.591 & 0.999 & 0.594 & 0.406 \\
budget 128 full & 761 & 0.329 & 0.897 & 0.568 & 0.997 & 0.651 & 0.413 \\
budget 128 full & 823 & 0.387 & 0.935 & 0.549 & 0.996 & 0.615 & 0.392 \\
budget 128 full & 887 & 0.377 & 0.914 & 0.537 & 0.998 & 0.609 & 0.360 \\
budget 128 full & 947 & 0.342 & 0.919 & 0.577 & 0.998 & 0.599 & 0.372 \\
obs full & 11 & 0.358 & 0.843 & 0.485 & 0.991 & 0.573 & 0.392 \\
obs full & 29 & 0.395 & 0.856 & 0.460 & 0.992 & 0.620 & 0.400 \\
obs full & 47 & 0.342 & 0.867 & 0.525 & 0.996 & 0.609 & 0.385 \\
obs full & 83 & 0.308 & 0.872 & 0.564 & 0.993 & 0.630 & 0.409 \\
obs full & 131 & 0.319 & 0.843 & 0.524 & 0.984 & 0.646 & 0.381 \\
obs full & 173 & 0.328 & 0.935 & 0.607 & 1.000 & 0.562 & 0.381 \\
obs full & 211 & 0.363 & 0.835 & 0.473 & 0.987 & 0.630 & 0.398 \\
obs full & 257 & 0.359 & 0.871 & 0.512 & 0.995 & 0.516 & 0.375 \\
obs full & 307 & 0.357 & 0.868 & 0.510 & 0.991 & 0.594 & 0.380 \\
obs full & 359 & 0.313 & 0.890 & 0.577 & 0.993 & 0.568 & 0.364 \\
obs full & 401 & 0.385 & 0.809 & 0.424 & 0.989 & 0.573 & 0.374 \\
obs full & 463 & 0.381 & 0.870 & 0.489 & 0.994 & 0.557 & 0.407 \\
obs full & 509 & 0.329 & 0.893 & 0.564 & 0.989 & 0.604 & 0.381 \\
obs full & 571 & 0.381 & 0.836 & 0.454 & 0.993 & 0.599 & 0.397 \\
obs full & 631 & 0.364 & 0.845 & 0.481 & 0.984 & 0.562 & 0.388 \\
obs full & 701 & 0.318 & 0.831 & 0.513 & 0.989 & 0.594 & 0.406 \\
obs full & 761 & 0.329 & 0.859 & 0.531 & 0.996 & 0.651 & 0.413 \\
obs full & 823 & 0.387 & 0.891 & 0.504 & 0.996 & 0.615 & 0.392 \\
obs full & 887 & 0.377 & 0.894 & 0.517 & 0.995 & 0.609 & 0.360 \\
obs full & 947 & 0.342 & 0.881 & 0.539 & 0.996 & 0.599 & 0.372 \\
obs coarse & 11 & 0.361 & 0.701 & 0.341 & 1.000 & 0.568 & 0.440 \\
obs coarse & 29 & 0.378 & 0.710 & 0.331 & 1.000 & 0.630 & 0.453 \\
obs coarse & 47 & 0.351 & 0.702 & 0.352 & 1.000 & 0.646 & 0.453 \\
obs coarse & 83 & 0.321 & 0.773 & 0.452 & 1.000 & 0.578 & 0.456 \\
obs coarse & 131 & 0.344 & 0.727 & 0.382 & 1.000 & 0.615 & 0.434 \\
obs coarse & 173 & 0.355 & 0.728 & 0.373 & 1.000 & 0.609 & 0.446 \\
obs coarse & 211 & 0.366 & 0.735 & 0.370 & 0.997 & 0.630 & 0.453 \\
obs coarse & 257 & 0.380 & 0.767 & 0.386 & 1.000 & 0.547 & 0.441 \\
obs coarse & 307 & 0.342 & 0.710 & 0.368 & 1.000 & 0.620 & 0.452 \\
obs coarse & 359 & 0.328 & 0.760 & 0.432 & 1.000 & 0.599 & 0.461 \\
obs coarse & 401 & 0.371 & 0.758 & 0.387 & 1.000 & 0.625 & 0.443 \\
obs coarse & 463 & 0.367 & 0.703 & 0.336 & 1.000 & 0.604 & 0.457 \\
obs coarse & 509 & 0.308 & 0.750 & 0.442 & 1.000 & 0.635 & 0.441 \\
obs coarse & 571 & 0.379 & 0.760 & 0.382 & 1.000 & 0.625 & 0.464 \\
obs coarse & 631 & 0.357 & 0.757 & 0.400 & 0.998 & 0.557 & 0.433 \\
obs coarse & 701 & 0.345 & 0.714 & 0.369 & 0.998 & 0.656 & 0.477 \\
obs coarse & 761 & 0.359 & 0.707 & 0.348 & 0.997 & 0.620 & 0.464 \\
obs coarse & 823 & 0.375 & 0.761 & 0.386 & 1.000 & 0.573 & 0.445 \\
obs coarse & 887 & 0.380 & 0.756 & 0.375 & 0.999 & 0.594 & 0.425 \\
obs coarse & 947 & 0.347 & 0.784 & 0.437 & 1.000 & 0.609 & 0.454 \\
obs blind & 11 & 0.358 & 0.400 & 0.042 & 1.000 & 0.568 & 0.442 \\
obs blind & 29 & 0.376 & 0.501 & 0.125 & 1.000 & 0.630 & 0.444 \\
obs blind & 47 & 0.342 & 0.423 & 0.081 & 1.000 & 0.651 & 0.454 \\
obs blind & 83 & 0.321 & 0.459 & 0.138 & 1.000 & 0.589 & 0.457 \\
obs blind & 131 & 0.340 & 0.463 & 0.124 & 1.000 & 0.615 & 0.426 \\
obs blind & 173 & 0.360 & 0.485 & 0.125 & 1.000 & 0.599 & 0.442 \\
obs blind & 211 & 0.355 & 0.519 & 0.164 & 1.000 & 0.615 & 0.450 \\
obs blind & 257 & 0.384 & 0.487 & 0.103 & 1.000 & 0.542 & 0.435 \\
obs blind & 307 & 0.336 & 0.417 & 0.081 & 1.000 & 0.615 & 0.441 \\
obs blind & 359 & 0.330 & 0.521 & 0.190 & 1.000 & 0.604 & 0.444 \\
obs blind & 401 & 0.360 & 0.456 & 0.096 & 1.000 & 0.620 & 0.437 \\
obs blind & 463 & 0.372 & 0.477 & 0.104 & 1.000 & 0.594 & 0.451 \\
obs blind & 509 & 0.297 & 0.486 & 0.189 & 1.000 & 0.641 & 0.436 \\
obs blind & 571 & 0.377 & 0.532 & 0.155 & 1.000 & 0.620 & 0.454 \\
obs blind & 631 & 0.349 & 0.491 & 0.143 & 1.000 & 0.562 & 0.426 \\
obs blind & 701 & 0.341 & 0.487 & 0.145 & 1.000 & 0.641 & 0.470 \\
obs blind & 761 & 0.352 & 0.486 & 0.134 & 1.000 & 0.630 & 0.468 \\
obs blind & 823 & 0.375 & 0.551 & 0.177 & 1.000 & 0.562 & 0.435 \\
obs blind & 887 & 0.374 & 0.500 & 0.126 & 1.000 & 0.604 & 0.423 \\
obs blind & 947 & 0.351 & 0.499 & 0.148 & 1.000 & 0.604 & 0.449 \\
blind undercovered & 11 & 0.518 & 0.460 & -0.058 & 0.694 & 0.398 & 0.313 \\
blind undercovered & 29 & 0.472 & 0.496 & 0.024 & 0.697 & 0.414 & 0.324 \\
blind undercovered & 47 & 0.518 & 0.450 & -0.068 & 0.683 & 0.434 & 0.330 \\
blind undercovered & 83 & 0.512 & 0.413 & -0.099 & 0.689 & 0.430 & 0.357 \\
blind undercovered & 131 & 0.525 & 0.464 & -0.061 & 0.681 & 0.434 & 0.315 \\
blind undercovered & 173 & 0.525 & 0.505 & -0.020 & 0.680 & 0.469 & 0.316 \\
blind undercovered & 211 & 0.508 & 0.494 & -0.015 & 0.670 & 0.461 & 0.327 \\
blind undercovered & 257 & 0.502 & 0.432 & -0.070 & 0.699 & 0.406 & 0.309 \\
blind undercovered & 307 & 0.536 & 0.478 & -0.057 & 0.697 & 0.477 & 0.339 \\
blind undercovered & 359 & 0.525 & 0.420 & -0.105 & 0.706 & 0.426 & 0.315 \\
blind undercovered & 401 & 0.502 & 0.408 & -0.094 & 0.704 & 0.398 & 0.334 \\
blind undercovered & 463 & 0.503 & 0.471 & -0.031 & 0.692 & 0.391 & 0.352 \\
blind undercovered & 509 & 0.525 & 0.495 & -0.030 & 0.705 & 0.434 & 0.348 \\
blind undercovered & 571 & 0.509 & 0.496 & -0.014 & 0.692 & 0.418 & 0.362 \\
blind undercovered & 631 & 0.473 & 0.432 & -0.040 & 0.683 & 0.418 & 0.297 \\
blind undercovered & 701 & 0.513 & 0.463 & -0.050 & 0.686 & 0.430 & 0.323 \\
blind undercovered & 761 & 0.498 & 0.401 & -0.098 & 0.698 & 0.438 & 0.316 \\
blind undercovered & 823 & 0.520 & 0.448 & -0.072 & 0.703 & 0.449 & 0.332 \\
blind undercovered & 887 & 0.496 & 0.447 & -0.049 & 0.706 & 0.426 & 0.343 \\
blind undercovered & 947 & 0.513 & 0.419 & -0.094 & 0.688 & 0.449 & 0.322 \\
sparse group04 & 11 & 0.282 & 0.832 & 0.551 & 0.994 & 0.486 & 0.407 \\
sparse group04 & 29 & 0.290 & 0.848 & 0.558 & 0.997 & 0.528 & 0.442 \\
sparse group04 & 47 & 0.301 & 0.829 & 0.528 & 0.991 & 0.472 & 0.413 \\
sparse group04 & 83 & 0.246 & 0.853 & 0.607 & 0.998 & 0.458 & 0.432 \\
sparse group04 & 131 & 0.320 & 0.838 & 0.518 & 0.996 & 0.479 & 0.433 \\
sparse group04 & 173 & 0.309 & 0.880 & 0.571 & 0.994 & 0.472 & 0.399 \\
sparse group04 & 211 & 0.286 & 0.852 & 0.566 & 0.993 & 0.493 & 0.404 \\
sparse group04 & 257 & 0.274 & 0.858 & 0.584 & 0.996 & 0.431 & 0.419 \\
sparse group04 & 307 & 0.303 & 0.857 & 0.555 & 0.991 & 0.403 & 0.412 \\
sparse group04 & 359 & 0.244 & 0.820 & 0.575 & 0.992 & 0.493 & 0.412 \\
sparse group04 & 401 & 0.278 & 0.817 & 0.539 & 0.988 & 0.479 & 0.417 \\
sparse group04 & 463 & 0.305 & 0.870 & 0.564 & 0.996 & 0.507 & 0.435 \\
sparse group04 & 509 & 0.296 & 0.864 & 0.568 & 0.986 & 0.424 & 0.439 \\
sparse group04 & 571 & 0.251 & 0.836 & 0.586 & 0.996 & 0.472 & 0.408 \\
sparse group04 & 631 & 0.287 & 0.845 & 0.558 & 0.993 & 0.451 & 0.399 \\
sparse group04 & 701 & 0.262 & 0.835 & 0.573 & 0.998 & 0.472 & 0.435 \\
sparse group04 & 761 & 0.286 & 0.781 & 0.495 & 0.987 & 0.493 & 0.429 \\
sparse group04 & 823 & 0.259 & 0.908 & 0.649 & 0.996 & 0.500 & 0.410 \\
sparse group04 & 887 & 0.287 & 0.861 & 0.574 & 0.994 & 0.479 & 0.412 \\
sparse group04 & 947 & 0.289 & 0.829 & 0.540 & 0.998 & 0.500 & 0.385 \\
sparse group08 & 11 & 0.337 & 0.924 & 0.587 & 0.999 & 0.497 & 0.409 \\
sparse group08 & 29 & 0.296 & 0.897 & 0.601 & 1.000 & 0.455 & 0.434 \\
sparse group08 & 47 & 0.324 & 0.894 & 0.570 & 0.995 & 0.462 & 0.425 \\
sparse group08 & 83 & 0.302 & 0.915 & 0.613 & 0.996 & 0.455 & 0.417 \\
sparse group08 & 131 & 0.307 & 0.931 & 0.624 & 0.998 & 0.458 & 0.430 \\
sparse group08 & 173 & 0.312 & 0.911 & 0.599 & 0.995 & 0.469 & 0.415 \\
sparse group08 & 211 & 0.299 & 0.921 & 0.622 & 0.999 & 0.455 & 0.405 \\
sparse group08 & 257 & 0.318 & 0.901 & 0.582 & 0.999 & 0.462 & 0.414 \\
sparse group08 & 307 & 0.302 & 0.912 & 0.610 & 0.998 & 0.417 & 0.419 \\
sparse group08 & 359 & 0.324 & 0.894 & 0.570 & 0.995 & 0.472 & 0.419 \\
sparse group08 & 401 & 0.296 & 0.885 & 0.589 & 0.996 & 0.479 & 0.417 \\
sparse group08 & 463 & 0.302 & 0.931 & 0.629 & 0.999 & 0.507 & 0.429 \\
sparse group08 & 509 & 0.318 & 0.923 & 0.605 & 0.994 & 0.476 & 0.428 \\
sparse group08 & 571 & 0.302 & 0.916 & 0.614 & 0.999 & 0.427 & 0.415 \\
sparse group08 & 631 & 0.321 & 0.904 & 0.583 & 0.997 & 0.465 & 0.417 \\
sparse group08 & 701 & 0.316 & 0.903 & 0.587 & 0.999 & 0.490 & 0.434 \\
sparse group08 & 761 & 0.298 & 0.914 & 0.616 & 0.997 & 0.479 & 0.434 \\
sparse group08 & 823 & 0.311 & 0.903 & 0.591 & 0.996 & 0.444 & 0.402 \\
sparse group08 & 887 & 0.306 & 0.916 & 0.610 & 0.997 & 0.476 & 0.419 \\
sparse group08 & 947 & 0.309 & 0.911 & 0.602 & 0.998 & 0.465 & 0.416 \\
\bottomrule
\end{longtable}


\normalsize

\section{Raw Error and Dispersion Rows}

Table \ref{tab:raw-error-rows} records the calibrated-MSE side of the same
matrix. These rows help distinguish two different claims: a method can have
higher correlation with oracle credit while also changing calibrated error,
zero-variance group frequency, or success-rate mix.

\scriptsize
% Generated by scripts/build_latex_inputs.py; do not edit by hand.
\begin{longtable}{@{}L{3.55cm}rrrrrrr@{}}
\caption{Raw case-by-seed calibrated-error rows from the deep matrix.}\label{tab:raw-error-rows}\\
\toprule
Case & Seed & Group MSE & Critic MSE & $\Delta$ MSE & Zero-var & Wait & Success \\
\midrule
\endfirsthead
\caption[]{Raw case-by-seed calibrated-error rows, continued.}\\
\toprule
Case & Seed & Group MSE & Critic MSE & $\Delta$ MSE & Zero-var & Wait & Success \\
\midrule
\endhead
\midrule
\multicolumn{8}{r}{\footnotesize\emph{Continued on next page}}\\
\endfoot
horizon 04 baseline & 11 & 0.050 & 0.004 & -0.046 & 0.000 & 0.290 & 0.422 \\
horizon 04 baseline & 29 & 0.053 & 0.003 & -0.050 & 0.000 & 0.312 & 0.391 \\
horizon 04 baseline & 47 & 0.047 & 0.002 & -0.045 & 0.000 & 0.286 & 0.396 \\
horizon 04 baseline & 83 & 0.044 & 0.003 & -0.041 & 0.000 & 0.331 & 0.458 \\
horizon 04 baseline & 131 & 0.047 & 0.003 & -0.044 & 0.000 & 0.297 & 0.391 \\
horizon 04 baseline & 173 & 0.049 & 0.004 & -0.046 & 0.000 & 0.296 & 0.422 \\
horizon 04 baseline & 211 & 0.053 & 0.002 & -0.051 & 0.000 & 0.310 & 0.422 \\
horizon 04 baseline & 257 & 0.054 & 0.003 & -0.051 & 0.000 & 0.305 & 0.427 \\
horizon 04 baseline & 307 & 0.044 & 0.002 & -0.042 & 0.000 & 0.317 & 0.438 \\
horizon 04 baseline & 359 & 0.058 & 0.003 & -0.055 & 0.000 & 0.295 & 0.422 \\
horizon 04 baseline & 401 & 0.051 & 0.003 & -0.048 & 0.000 & 0.298 & 0.438 \\
horizon 04 baseline & 463 & 0.051 & 0.003 & -0.048 & 0.000 & 0.309 & 0.417 \\
horizon 04 baseline & 509 & 0.054 & 0.002 & -0.052 & 0.000 & 0.294 & 0.385 \\
horizon 04 baseline & 571 & 0.052 & 0.004 & -0.048 & 0.000 & 0.330 & 0.396 \\
horizon 04 baseline & 631 & 0.056 & 0.003 & -0.053 & 0.000 & 0.265 & 0.474 \\
horizon 04 baseline & 701 & 0.052 & 0.004 & -0.049 & 0.000 & 0.304 & 0.448 \\
horizon 04 baseline & 761 & 0.047 & 0.007 & -0.040 & 0.000 & 0.299 & 0.443 \\
horizon 04 baseline & 823 & 0.045 & 0.003 & -0.042 & 0.000 & 0.312 & 0.464 \\
horizon 04 baseline & 887 & 0.049 & 0.002 & -0.047 & 0.000 & 0.301 & 0.422 \\
horizon 04 baseline & 947 & 0.050 & 0.002 & -0.048 & 0.000 & 0.304 & 0.448 \\
\addlinespace[2pt]
horizon 08 baseline & 11 & 0.034 & 0.008 & -0.026 & 0.000 & 0.361 & 0.542 \\
horizon 08 baseline & 29 & 0.036 & 0.007 & -0.029 & 0.000 & 0.382 & 0.531 \\
horizon 08 baseline & 47 & 0.036 & 0.008 & -0.028 & 0.000 & 0.395 & 0.542 \\
horizon 08 baseline & 83 & 0.036 & 0.006 & -0.030 & 0.000 & 0.369 & 0.490 \\
horizon 08 baseline & 131 & 0.030 & 0.007 & -0.022 & 0.000 & 0.362 & 0.552 \\
horizon 08 baseline & 173 & 0.038 & 0.007 & -0.031 & 0.000 & 0.373 & 0.542 \\
horizon 08 baseline & 211 & 0.038 & 0.007 & -0.031 & 0.000 & 0.347 & 0.547 \\
horizon 08 baseline & 257 & 0.035 & 0.006 & -0.029 & 0.000 & 0.373 & 0.552 \\
horizon 08 baseline & 307 & 0.034 & 0.004 & -0.030 & 0.000 & 0.364 & 0.583 \\
horizon 08 baseline & 359 & 0.039 & 0.004 & -0.035 & 0.000 & 0.376 & 0.526 \\
horizon 08 baseline & 401 & 0.035 & 0.008 & -0.027 & 0.000 & 0.334 & 0.516 \\
horizon 08 baseline & 463 & 0.037 & 0.006 & -0.031 & 0.000 & 0.380 & 0.589 \\
horizon 08 baseline & 509 & 0.034 & 0.006 & -0.028 & 0.000 & 0.369 & 0.599 \\
horizon 08 baseline & 571 & 0.040 & 0.003 & -0.037 & 0.000 & 0.347 & 0.568 \\
horizon 08 baseline & 631 & 0.035 & 0.004 & -0.031 & 0.000 & 0.341 & 0.510 \\
horizon 08 baseline & 701 & 0.033 & 0.006 & -0.027 & 0.000 & 0.404 & 0.583 \\
horizon 08 baseline & 761 & 0.037 & 0.007 & -0.030 & 0.000 & 0.350 & 0.547 \\
horizon 08 baseline & 823 & 0.031 & 0.006 & -0.025 & 0.000 & 0.357 & 0.578 \\
horizon 08 baseline & 887 & 0.032 & 0.005 & -0.028 & 0.000 & 0.364 & 0.604 \\
horizon 08 baseline & 947 & 0.037 & 0.005 & -0.032 & 0.000 & 0.358 & 0.500 \\
\addlinespace[2pt]
horizon 12 baseline & 11 & 0.024 & 0.007 & -0.017 & 0.000 & 0.415 & 0.609 \\
horizon 12 baseline & 29 & 0.028 & 0.005 & -0.023 & 0.000 & 0.400 & 0.589 \\
horizon 12 baseline & 47 & 0.022 & 0.006 & -0.016 & 0.000 & 0.417 & 0.677 \\
horizon 12 baseline & 83 & 0.027 & 0.007 & -0.019 & 0.000 & 0.396 & 0.573 \\
horizon 12 baseline & 131 & 0.026 & 0.007 & -0.019 & 0.000 & 0.397 & 0.604 \\
horizon 12 baseline & 173 & 0.026 & 0.009 & -0.017 & 0.000 & 0.416 & 0.630 \\
horizon 12 baseline & 211 & 0.025 & 0.007 & -0.018 & 0.000 & 0.390 & 0.667 \\
horizon 12 baseline & 257 & 0.023 & 0.005 & -0.017 & 0.000 & 0.437 & 0.656 \\
horizon 12 baseline & 307 & 0.025 & 0.006 & -0.018 & 0.000 & 0.412 & 0.667 \\
horizon 12 baseline & 359 & 0.024 & 0.008 & -0.016 & 0.000 & 0.397 & 0.599 \\
horizon 12 baseline & 401 & 0.025 & 0.007 & -0.018 & 0.000 & 0.412 & 0.630 \\
horizon 12 baseline & 463 & 0.022 & 0.007 & -0.015 & 0.000 & 0.416 & 0.656 \\
horizon 12 baseline & 509 & 0.025 & 0.006 & -0.019 & 0.000 & 0.425 & 0.661 \\
horizon 12 baseline & 571 & 0.024 & 0.008 & -0.016 & 0.000 & 0.424 & 0.661 \\
horizon 12 baseline & 631 & 0.028 & 0.009 & -0.019 & 0.000 & 0.392 & 0.604 \\
horizon 12 baseline & 701 & 0.024 & 0.005 & -0.019 & 0.000 & 0.445 & 0.651 \\
horizon 12 baseline & 761 & 0.025 & 0.006 & -0.019 & 0.000 & 0.405 & 0.625 \\
horizon 12 baseline & 823 & 0.027 & 0.009 & -0.018 & 0.000 & 0.399 & 0.573 \\
horizon 12 baseline & 887 & 0.025 & 0.008 & -0.017 & 0.000 & 0.402 & 0.646 \\
horizon 12 baseline & 947 & 0.025 & 0.008 & -0.017 & 0.000 & 0.421 & 0.620 \\
\addlinespace[2pt]
horizon 16 long wait & 11 & 0.021 & 0.006 & -0.015 & 0.000 & 0.513 & 0.667 \\
horizon 16 long wait & 29 & 0.018 & 0.004 & -0.014 & 0.000 & 0.536 & 0.667 \\
horizon 16 long wait & 47 & 0.017 & 0.005 & -0.013 & 0.000 & 0.538 & 0.718 \\
horizon 16 long wait & 83 & 0.020 & 0.005 & -0.016 & 0.000 & 0.529 & 0.620 \\
horizon 16 long wait & 131 & 0.019 & 0.005 & -0.014 & 0.000 & 0.542 & 0.671 \\
horizon 16 long wait & 173 & 0.019 & 0.005 & -0.014 & 0.000 & 0.552 & 0.606 \\
horizon 16 long wait & 211 & 0.019 & 0.005 & -0.014 & 0.000 & 0.546 & 0.718 \\
horizon 16 long wait & 257 & 0.017 & 0.005 & -0.012 & 0.000 & 0.532 & 0.694 \\
horizon 16 long wait & 307 & 0.023 & 0.004 & -0.019 & 0.000 & 0.539 & 0.671 \\
horizon 16 long wait & 359 & 0.016 & 0.005 & -0.011 & 0.000 & 0.546 & 0.741 \\
horizon 16 long wait & 401 & 0.017 & 0.006 & -0.011 & 0.000 & 0.539 & 0.718 \\
horizon 16 long wait & 463 & 0.018 & 0.005 & -0.013 & 0.000 & 0.552 & 0.662 \\
horizon 16 long wait & 509 & 0.019 & 0.005 & -0.014 & 0.000 & 0.545 & 0.662 \\
horizon 16 long wait & 571 & 0.018 & 0.006 & -0.012 & 0.000 & 0.547 & 0.657 \\
horizon 16 long wait & 631 & 0.019 & 0.006 & -0.013 & 0.000 & 0.517 & 0.667 \\
horizon 16 long wait & 701 & 0.017 & 0.005 & -0.012 & 0.000 & 0.527 & 0.713 \\
horizon 16 long wait & 761 & 0.015 & 0.004 & -0.011 & 0.000 & 0.562 & 0.736 \\
horizon 16 long wait & 823 & 0.022 & 0.006 & -0.016 & 0.000 & 0.532 & 0.657 \\
horizon 16 long wait & 887 & 0.020 & 0.004 & -0.016 & 0.000 & 0.539 & 0.713 \\
horizon 16 long wait & 947 & 0.018 & 0.005 & -0.013 & 0.000 & 0.545 & 0.694 \\
\addlinespace[2pt]
group 02 long wait & 11 & 0.024 & 0.014 & -0.010 & 0.031 & 0.507 & 0.641 \\
group 02 long wait & 29 & 0.024 & 0.010 & -0.014 & 0.094 & 0.496 & 0.719 \\
group 02 long wait & 47 & 0.025 & 0.010 & -0.015 & 0.062 & 0.495 & 0.656 \\
group 02 long wait & 83 & 0.022 & 0.011 & -0.010 & 0.125 & 0.487 & 0.719 \\
group 02 long wait & 131 & 0.024 & 0.010 & -0.013 & 0.094 & 0.514 & 0.625 \\
group 02 long wait & 173 & 0.024 & 0.009 & -0.015 & 0.125 & 0.499 & 0.734 \\
group 02 long wait & 211 & 0.025 & 0.010 & -0.014 & 0.031 & 0.489 & 0.578 \\
group 02 long wait & 257 & 0.024 & 0.010 & -0.014 & 0.156 & 0.515 & 0.641 \\
group 02 long wait & 307 & 0.026 & 0.009 & -0.017 & 0.062 & 0.533 & 0.672 \\
group 02 long wait & 359 & 0.026 & 0.014 & -0.012 & 0.000 & 0.453 & 0.625 \\
group 02 long wait & 401 & 0.024 & 0.012 & -0.012 & 0.062 & 0.528 & 0.672 \\
group 02 long wait & 463 & 0.022 & 0.008 & -0.014 & 0.062 & 0.496 & 0.609 \\
group 02 long wait & 509 & 0.028 & 0.013 & -0.015 & 0.125 & 0.515 & 0.656 \\
group 02 long wait & 571 & 0.035 & 0.009 & -0.026 & 0.062 & 0.536 & 0.656 \\
group 02 long wait & 631 & 0.027 & 0.011 & -0.016 & 0.000 & 0.492 & 0.641 \\
group 02 long wait & 701 & 0.028 & 0.013 & -0.015 & 0.062 & 0.488 & 0.609 \\
group 02 long wait & 761 & 0.027 & 0.011 & -0.016 & 0.031 & 0.490 & 0.672 \\
group 02 long wait & 823 & 0.024 & 0.009 & -0.015 & 0.062 & 0.498 & 0.688 \\
group 02 long wait & 887 & 0.017 & 0.007 & -0.010 & 0.031 & 0.537 & 0.734 \\
group 02 long wait & 947 & 0.024 & 0.008 & -0.016 & 0.031 & 0.542 & 0.703 \\
\addlinespace[2pt]
group 04 long wait & 11 & 0.023 & 0.010 & -0.014 & 0.000 & 0.488 & 0.680 \\
group 04 long wait & 29 & 0.022 & 0.006 & -0.016 & 0.000 & 0.495 & 0.734 \\
group 04 long wait & 47 & 0.023 & 0.009 & -0.014 & 0.000 & 0.481 & 0.656 \\
group 04 long wait & 83 & 0.022 & 0.006 & -0.016 & 0.000 & 0.503 & 0.664 \\
group 04 long wait & 131 & 0.025 & 0.007 & -0.018 & 0.000 & 0.524 & 0.625 \\
group 04 long wait & 173 & 0.025 & 0.010 & -0.015 & 0.000 & 0.465 & 0.656 \\
group 04 long wait & 211 & 0.026 & 0.008 & -0.018 & 0.000 & 0.507 & 0.609 \\
group 04 long wait & 257 & 0.021 & 0.006 & -0.016 & 0.000 & 0.533 & 0.664 \\
group 04 long wait & 307 & 0.025 & 0.006 & -0.019 & 0.000 & 0.495 & 0.641 \\
group 04 long wait & 359 & 0.029 & 0.008 & -0.021 & 0.000 & 0.471 & 0.609 \\
group 04 long wait & 401 & 0.022 & 0.004 & -0.018 & 0.000 & 0.544 & 0.656 \\
group 04 long wait & 463 & 0.027 & 0.007 & -0.021 & 0.000 & 0.487 & 0.570 \\
group 04 long wait & 509 & 0.028 & 0.010 & -0.018 & 0.000 & 0.516 & 0.617 \\
group 04 long wait & 571 & 0.026 & 0.007 & -0.019 & 0.000 & 0.514 & 0.688 \\
group 04 long wait & 631 & 0.021 & 0.007 & -0.014 & 0.000 & 0.485 & 0.680 \\
group 04 long wait & 701 & 0.023 & 0.006 & -0.017 & 0.000 & 0.517 & 0.688 \\
group 04 long wait & 761 & 0.022 & 0.006 & -0.017 & 0.000 & 0.491 & 0.688 \\
group 04 long wait & 823 & 0.024 & 0.007 & -0.017 & 0.000 & 0.508 & 0.664 \\
group 04 long wait & 887 & 0.024 & 0.009 & -0.015 & 0.000 & 0.508 & 0.672 \\
group 04 long wait & 947 & 0.020 & 0.007 & -0.014 & 0.000 & 0.516 & 0.633 \\
\addlinespace[2pt]
group 08 long wait & 11 & 0.024 & 0.005 & -0.019 & 0.000 & 0.506 & 0.656 \\
group 08 long wait & 29 & 0.026 & 0.005 & -0.021 & 0.000 & 0.493 & 0.652 \\
group 08 long wait & 47 & 0.022 & 0.004 & -0.018 & 0.000 & 0.506 & 0.680 \\
group 08 long wait & 83 & 0.025 & 0.004 & -0.021 & 0.000 & 0.460 & 0.625 \\
group 08 long wait & 131 & 0.022 & 0.003 & -0.019 & 0.000 & 0.522 & 0.637 \\
group 08 long wait & 173 & 0.025 & 0.007 & -0.018 & 0.000 & 0.497 & 0.633 \\
group 08 long wait & 211 & 0.026 & 0.004 & -0.021 & 0.000 & 0.503 & 0.621 \\
group 08 long wait & 257 & 0.025 & 0.003 & -0.021 & 0.000 & 0.505 & 0.617 \\
group 08 long wait & 307 & 0.026 & 0.004 & -0.021 & 0.000 & 0.511 & 0.664 \\
group 08 long wait & 359 & 0.024 & 0.006 & -0.019 & 0.000 & 0.505 & 0.660 \\
group 08 long wait & 401 & 0.024 & 0.004 & -0.020 & 0.000 & 0.515 & 0.660 \\
group 08 long wait & 463 & 0.025 & 0.005 & -0.020 & 0.000 & 0.518 & 0.609 \\
group 08 long wait & 509 & 0.024 & 0.004 & -0.020 & 0.000 & 0.503 & 0.691 \\
group 08 long wait & 571 & 0.023 & 0.007 & -0.017 & 0.000 & 0.526 & 0.672 \\
group 08 long wait & 631 & 0.021 & 0.004 & -0.017 & 0.000 & 0.517 & 0.609 \\
group 08 long wait & 701 & 0.021 & 0.005 & -0.017 & 0.000 & 0.505 & 0.668 \\
group 08 long wait & 761 & 0.026 & 0.003 & -0.023 & 0.000 & 0.494 & 0.586 \\
group 08 long wait & 823 & 0.024 & 0.006 & -0.018 & 0.000 & 0.484 & 0.617 \\
group 08 long wait & 887 & 0.026 & 0.005 & -0.021 & 0.000 & 0.484 & 0.613 \\
group 08 long wait & 947 & 0.020 & 0.004 & -0.016 & 0.000 & 0.513 & 0.660 \\
\addlinespace[2pt]
group 12 long wait & 11 & 0.024 & 0.004 & -0.020 & 0.000 & 0.496 & 0.646 \\
group 12 long wait & 29 & 0.026 & 0.004 & -0.023 & 0.000 & 0.491 & 0.633 \\
group 12 long wait & 47 & 0.023 & 0.002 & -0.021 & 0.000 & 0.511 & 0.659 \\
group 12 long wait & 83 & 0.023 & 0.003 & -0.020 & 0.000 & 0.498 & 0.674 \\
group 12 long wait & 131 & 0.024 & 0.002 & -0.022 & 0.000 & 0.531 & 0.638 \\
group 12 long wait & 173 & 0.023 & 0.003 & -0.021 & 0.000 & 0.490 & 0.690 \\
group 12 long wait & 211 & 0.028 & 0.003 & -0.024 & 0.000 & 0.488 & 0.617 \\
group 12 long wait & 257 & 0.025 & 0.003 & -0.021 & 0.000 & 0.511 & 0.607 \\
group 12 long wait & 307 & 0.025 & 0.003 & -0.022 & 0.000 & 0.502 & 0.654 \\
group 12 long wait & 359 & 0.024 & 0.003 & -0.021 & 0.000 & 0.493 & 0.667 \\
group 12 long wait & 401 & 0.023 & 0.003 & -0.020 & 0.000 & 0.502 & 0.674 \\
group 12 long wait & 463 & 0.024 & 0.004 & -0.020 & 0.000 & 0.520 & 0.646 \\
group 12 long wait & 509 & 0.025 & 0.003 & -0.022 & 0.000 & 0.490 & 0.635 \\
group 12 long wait & 571 & 0.024 & 0.004 & -0.020 & 0.000 & 0.511 & 0.628 \\
group 12 long wait & 631 & 0.021 & 0.003 & -0.018 & 0.000 & 0.507 & 0.656 \\
group 12 long wait & 701 & 0.021 & 0.003 & -0.018 & 0.000 & 0.513 & 0.693 \\
group 12 long wait & 761 & 0.024 & 0.003 & -0.020 & 0.000 & 0.500 & 0.659 \\
group 12 long wait & 823 & 0.026 & 0.003 & -0.023 & 0.000 & 0.486 & 0.617 \\
group 12 long wait & 887 & 0.026 & 0.003 & -0.023 & 0.000 & 0.496 & 0.635 \\
group 12 long wait & 947 & 0.023 & 0.003 & -0.020 & 0.000 & 0.494 & 0.656 \\
\addlinespace[2pt]
budget 002 full & 11 & 0.029 & 0.027 & -0.001 & 0.000 & 0.392 & 0.573 \\
budget 002 full & 29 & 0.028 & 0.028 & -0.000 & 0.000 & 0.400 & 0.620 \\
budget 002 full & 47 & 0.029 & 0.029 & 0.000 & 0.000 & 0.385 & 0.609 \\
budget 002 full & 83 & 0.027 & 0.028 & 0.001 & 0.000 & 0.409 & 0.630 \\
budget 002 full & 131 & 0.028 & 0.028 & -0.000 & 0.000 & 0.381 & 0.646 \\
budget 002 full & 173 & 0.036 & 0.033 & -0.004 & 0.000 & 0.381 & 0.562 \\
budget 002 full & 211 & 0.028 & 0.029 & 0.001 & 0.000 & 0.398 & 0.630 \\
budget 002 full & 257 & 0.034 & 0.031 & -0.003 & 0.000 & 0.375 & 0.516 \\
budget 002 full & 307 & 0.029 & 0.031 & 0.002 & 0.000 & 0.380 & 0.594 \\
budget 002 full & 359 & 0.035 & 0.032 & -0.003 & 0.000 & 0.364 & 0.568 \\
budget 002 full & 401 & 0.030 & 0.032 & 0.002 & 0.000 & 0.374 & 0.573 \\
budget 002 full & 463 & 0.033 & 0.029 & -0.003 & 0.000 & 0.407 & 0.557 \\
budget 002 full & 509 & 0.030 & 0.030 & -0.001 & 0.000 & 0.381 & 0.604 \\
budget 002 full & 571 & 0.031 & 0.031 & 0.000 & 0.000 & 0.397 & 0.599 \\
budget 002 full & 631 & 0.031 & 0.032 & 0.001 & 0.000 & 0.388 & 0.562 \\
budget 002 full & 701 & 0.032 & 0.031 & -0.001 & 0.000 & 0.406 & 0.594 \\
budget 002 full & 761 & 0.026 & 0.025 & -0.001 & 0.000 & 0.413 & 0.651 \\
budget 002 full & 823 & 0.027 & 0.026 & -0.001 & 0.000 & 0.392 & 0.615 \\
budget 002 full & 887 & 0.031 & 0.032 & 0.000 & 0.000 & 0.360 & 0.609 \\
budget 002 full & 947 & 0.033 & 0.032 & -0.001 & 0.000 & 0.372 & 0.599 \\
\addlinespace[2pt]
budget 008 full & 11 & 0.029 & 0.026 & -0.003 & 0.000 & 0.392 & 0.573 \\
budget 008 full & 29 & 0.028 & 0.024 & -0.004 & 0.000 & 0.400 & 0.620 \\
budget 008 full & 47 & 0.029 & 0.024 & -0.005 & 0.000 & 0.385 & 0.609 \\
budget 008 full & 83 & 0.027 & 0.024 & -0.004 & 0.000 & 0.409 & 0.630 \\
budget 008 full & 131 & 0.028 & 0.026 & -0.002 & 0.000 & 0.381 & 0.646 \\
budget 008 full & 173 & 0.036 & 0.027 & -0.009 & 0.000 & 0.381 & 0.562 \\
budget 008 full & 211 & 0.028 & 0.025 & -0.003 & 0.000 & 0.398 & 0.630 \\
budget 008 full & 257 & 0.034 & 0.028 & -0.005 & 0.000 & 0.375 & 0.516 \\
budget 008 full & 307 & 0.029 & 0.028 & -0.001 & 0.000 & 0.380 & 0.594 \\
budget 008 full & 359 & 0.035 & 0.028 & -0.007 & 0.000 & 0.364 & 0.568 \\
budget 008 full & 401 & 0.030 & 0.027 & -0.003 & 0.000 & 0.374 & 0.573 \\
budget 008 full & 463 & 0.033 & 0.030 & -0.002 & 0.000 & 0.407 & 0.557 \\
budget 008 full & 509 & 0.030 & 0.028 & -0.003 & 0.000 & 0.381 & 0.604 \\
budget 008 full & 571 & 0.031 & 0.028 & -0.003 & 0.000 & 0.397 & 0.599 \\
budget 008 full & 631 & 0.031 & 0.031 & 0.000 & 0.000 & 0.388 & 0.562 \\
budget 008 full & 701 & 0.032 & 0.029 & -0.003 & 0.000 & 0.406 & 0.594 \\
budget 008 full & 761 & 0.026 & 0.026 & -0.000 & 0.000 & 0.413 & 0.651 \\
budget 008 full & 823 & 0.027 & 0.026 & -0.001 & 0.000 & 0.392 & 0.615 \\
budget 008 full & 887 & 0.031 & 0.027 & -0.004 & 0.000 & 0.360 & 0.609 \\
budget 008 full & 947 & 0.033 & 0.030 & -0.003 & 0.000 & 0.372 & 0.599 \\
\addlinespace[2pt]
budget 032 full & 11 & 0.029 & 0.016 & -0.013 & 0.000 & 0.392 & 0.573 \\
budget 032 full & 29 & 0.028 & 0.014 & -0.014 & 0.000 & 0.400 & 0.620 \\
budget 032 full & 47 & 0.029 & 0.014 & -0.014 & 0.000 & 0.385 & 0.609 \\
budget 032 full & 83 & 0.027 & 0.015 & -0.012 & 0.000 & 0.409 & 0.630 \\
budget 032 full & 131 & 0.028 & 0.015 & -0.013 & 0.000 & 0.381 & 0.646 \\
budget 032 full & 173 & 0.036 & 0.013 & -0.023 & 0.000 & 0.381 & 0.562 \\
budget 032 full & 211 & 0.028 & 0.018 & -0.010 & 0.000 & 0.398 & 0.630 \\
budget 032 full & 257 & 0.034 & 0.014 & -0.019 & 0.000 & 0.375 & 0.516 \\
budget 032 full & 307 & 0.029 & 0.017 & -0.013 & 0.000 & 0.380 & 0.594 \\
budget 032 full & 359 & 0.035 & 0.015 & -0.020 & 0.000 & 0.364 & 0.568 \\
budget 032 full & 401 & 0.030 & 0.018 & -0.012 & 0.000 & 0.374 & 0.573 \\
budget 032 full & 463 & 0.033 & 0.016 & -0.016 & 0.000 & 0.407 & 0.557 \\
budget 032 full & 509 & 0.030 & 0.019 & -0.012 & 0.000 & 0.381 & 0.604 \\
budget 032 full & 571 & 0.031 & 0.018 & -0.013 & 0.000 & 0.397 & 0.599 \\
budget 032 full & 631 & 0.031 & 0.017 & -0.014 & 0.000 & 0.388 & 0.562 \\
budget 032 full & 701 & 0.032 & 0.018 & -0.014 & 0.000 & 0.406 & 0.594 \\
budget 032 full & 761 & 0.026 & 0.015 & -0.011 & 0.000 & 0.413 & 0.651 \\
budget 032 full & 823 & 0.027 & 0.011 & -0.015 & 0.000 & 0.392 & 0.615 \\
budget 032 full & 887 & 0.031 & 0.015 & -0.017 & 0.000 & 0.360 & 0.609 \\
budget 032 full & 947 & 0.033 & 0.017 & -0.016 & 0.000 & 0.372 & 0.599 \\
\addlinespace[2pt]
budget 128 full & 11 & 0.029 & 0.006 & -0.023 & 0.000 & 0.392 & 0.573 \\
budget 128 full & 29 & 0.028 & 0.007 & -0.021 & 0.000 & 0.400 & 0.620 \\
budget 128 full & 47 & 0.029 & 0.006 & -0.023 & 0.000 & 0.385 & 0.609 \\
budget 128 full & 83 & 0.027 & 0.006 & -0.021 & 0.000 & 0.409 & 0.630 \\
budget 128 full & 131 & 0.028 & 0.006 & -0.022 & 0.000 & 0.381 & 0.646 \\
budget 128 full & 173 & 0.036 & 0.005 & -0.032 & 0.000 & 0.381 & 0.562 \\
budget 128 full & 211 & 0.028 & 0.008 & -0.020 & 0.000 & 0.398 & 0.630 \\
budget 128 full & 257 & 0.034 & 0.006 & -0.027 & 0.000 & 0.375 & 0.516 \\
budget 128 full & 307 & 0.029 & 0.008 & -0.022 & 0.000 & 0.380 & 0.594 \\
budget 128 full & 359 & 0.035 & 0.005 & -0.030 & 0.000 & 0.364 & 0.568 \\
budget 128 full & 401 & 0.030 & 0.006 & -0.024 & 0.000 & 0.374 & 0.573 \\
budget 128 full & 463 & 0.033 & 0.006 & -0.027 & 0.000 & 0.407 & 0.557 \\
budget 128 full & 509 & 0.030 & 0.005 & -0.025 & 0.000 & 0.381 & 0.604 \\
budget 128 full & 571 & 0.031 & 0.008 & -0.023 & 0.000 & 0.397 & 0.599 \\
budget 128 full & 631 & 0.031 & 0.008 & -0.023 & 0.000 & 0.388 & 0.562 \\
budget 128 full & 701 & 0.032 & 0.006 & -0.026 & 0.000 & 0.406 & 0.594 \\
budget 128 full & 761 & 0.026 & 0.006 & -0.020 & 0.000 & 0.413 & 0.651 \\
budget 128 full & 823 & 0.027 & 0.004 & -0.023 & 0.000 & 0.392 & 0.615 \\
budget 128 full & 887 & 0.031 & 0.006 & -0.025 & 0.000 & 0.360 & 0.609 \\
budget 128 full & 947 & 0.033 & 0.006 & -0.027 & 0.000 & 0.372 & 0.599 \\
\addlinespace[2pt]
obs full & 11 & 0.029 & 0.010 & -0.019 & 0.000 & 0.392 & 0.573 \\
obs full & 29 & 0.028 & 0.009 & -0.019 & 0.000 & 0.400 & 0.620 \\
obs full & 47 & 0.029 & 0.008 & -0.021 & 0.000 & 0.385 & 0.609 \\
obs full & 83 & 0.027 & 0.007 & -0.020 & 0.000 & 0.409 & 0.630 \\
obs full & 131 & 0.028 & 0.009 & -0.019 & 0.000 & 0.381 & 0.646 \\
obs full & 173 & 0.036 & 0.005 & -0.031 & 0.000 & 0.381 & 0.562 \\
obs full & 211 & 0.028 & 0.010 & -0.018 & 0.000 & 0.398 & 0.630 \\
obs full & 257 & 0.034 & 0.009 & -0.024 & 0.000 & 0.375 & 0.516 \\
obs full & 307 & 0.029 & 0.008 & -0.021 & 0.000 & 0.380 & 0.594 \\
obs full & 359 & 0.035 & 0.008 & -0.027 & 0.000 & 0.364 & 0.568 \\
obs full & 401 & 0.030 & 0.012 & -0.018 & 0.000 & 0.374 & 0.573 \\
obs full & 463 & 0.033 & 0.009 & -0.024 & 0.000 & 0.407 & 0.557 \\
obs full & 509 & 0.030 & 0.007 & -0.024 & 0.000 & 0.381 & 0.604 \\
obs full & 571 & 0.031 & 0.011 & -0.020 & 0.000 & 0.397 & 0.599 \\
obs full & 631 & 0.031 & 0.010 & -0.021 & 0.000 & 0.388 & 0.562 \\
obs full & 701 & 0.032 & 0.011 & -0.021 & 0.000 & 0.406 & 0.594 \\
obs full & 761 & 0.026 & 0.008 & -0.018 & 0.000 & 0.413 & 0.651 \\
obs full & 823 & 0.027 & 0.007 & -0.020 & 0.000 & 0.392 & 0.615 \\
obs full & 887 & 0.031 & 0.007 & -0.024 & 0.000 & 0.360 & 0.609 \\
obs full & 947 & 0.033 & 0.008 & -0.024 & 0.000 & 0.372 & 0.599 \\
\addlinespace[2pt]
obs coarse & 11 & 0.029 & 0.017 & -0.012 & 0.000 & 0.440 & 0.568 \\
obs coarse & 29 & 0.029 & 0.017 & -0.012 & 0.000 & 0.453 & 0.630 \\
obs coarse & 47 & 0.025 & 0.014 & -0.010 & 0.000 & 0.453 & 0.646 \\
obs coarse & 83 & 0.031 & 0.014 & -0.017 & 0.000 & 0.456 & 0.578 \\
obs coarse & 131 & 0.031 & 0.017 & -0.014 & 0.000 & 0.434 & 0.615 \\
obs coarse & 173 & 0.031 & 0.016 & -0.014 & 0.000 & 0.446 & 0.609 \\
obs coarse & 211 & 0.028 & 0.015 & -0.013 & 0.000 & 0.453 & 0.630 \\
obs coarse & 257 & 0.033 & 0.016 & -0.017 & 0.000 & 0.441 & 0.547 \\
obs coarse & 307 & 0.028 & 0.016 & -0.012 & 0.000 & 0.452 & 0.620 \\
obs coarse & 359 & 0.031 & 0.015 & -0.016 & 0.000 & 0.461 & 0.599 \\
obs coarse & 401 & 0.029 & 0.014 & -0.015 & 0.000 & 0.443 & 0.625 \\
obs coarse & 463 & 0.029 & 0.017 & -0.012 & 0.000 & 0.457 & 0.604 \\
obs coarse & 509 & 0.030 & 0.015 & -0.016 & 0.000 & 0.441 & 0.635 \\
obs coarse & 571 & 0.029 & 0.014 & -0.015 & 0.000 & 0.464 & 0.625 \\
obs coarse & 631 & 0.033 & 0.016 & -0.017 & 0.000 & 0.433 & 0.557 \\
obs coarse & 701 & 0.026 & 0.015 & -0.012 & 0.000 & 0.477 & 0.656 \\
obs coarse & 761 & 0.028 & 0.016 & -0.012 & 0.000 & 0.464 & 0.620 \\
obs coarse & 823 & 0.032 & 0.016 & -0.016 & 0.000 & 0.445 & 0.573 \\
obs coarse & 887 & 0.033 & 0.017 & -0.017 & 0.000 & 0.425 & 0.594 \\
obs coarse & 947 & 0.032 & 0.014 & -0.018 & 0.000 & 0.454 & 0.609 \\
\addlinespace[2pt]
obs blind & 11 & 0.029 & 0.028 & -0.001 & 0.000 & 0.442 & 0.568 \\
obs blind & 29 & 0.029 & 0.025 & -0.004 & 0.000 & 0.444 & 0.630 \\
obs blind & 47 & 0.025 & 0.023 & -0.002 & 0.000 & 0.454 & 0.651 \\
obs blind & 83 & 0.030 & 0.026 & -0.004 & 0.000 & 0.457 & 0.589 \\
obs blind & 131 & 0.031 & 0.027 & -0.003 & 0.000 & 0.426 & 0.615 \\
obs blind & 173 & 0.032 & 0.028 & -0.004 & 0.000 & 0.442 & 0.599 \\
obs blind & 211 & 0.029 & 0.025 & -0.005 & 0.000 & 0.450 & 0.615 \\
obs blind & 257 & 0.033 & 0.029 & -0.003 & 0.000 & 0.435 & 0.542 \\
obs blind & 307 & 0.029 & 0.027 & -0.002 & 0.000 & 0.441 & 0.615 \\
obs blind & 359 & 0.031 & 0.025 & -0.006 & 0.000 & 0.444 & 0.604 \\
obs blind & 401 & 0.030 & 0.027 & -0.003 & 0.000 & 0.437 & 0.620 \\
obs blind & 463 & 0.029 & 0.026 & -0.003 & 0.000 & 0.451 & 0.594 \\
obs blind & 509 & 0.030 & 0.025 & -0.005 & 0.000 & 0.436 & 0.641 \\
obs blind & 571 & 0.030 & 0.025 & -0.005 & 0.000 & 0.454 & 0.620 \\
obs blind & 631 & 0.033 & 0.029 & -0.005 & 0.000 & 0.426 & 0.562 \\
obs blind & 701 & 0.027 & 0.024 & -0.004 & 0.000 & 0.470 & 0.641 \\
obs blind & 761 & 0.028 & 0.024 & -0.004 & 0.000 & 0.468 & 0.630 \\
obs blind & 823 & 0.032 & 0.026 & -0.006 & 0.000 & 0.435 & 0.562 \\
obs blind & 887 & 0.032 & 0.028 & -0.004 & 0.000 & 0.423 & 0.604 \\
obs blind & 947 & 0.031 & 0.027 & -0.004 & 0.000 & 0.449 & 0.604 \\
\addlinespace[2pt]
blind undercov. & 11 & 0.051 & 0.055 & 0.004 & 0.000 & 0.313 & 0.398 \\
blind undercov. & 29 & 0.053 & 0.052 & -0.002 & 0.000 & 0.324 & 0.414 \\
blind undercov. & 47 & 0.048 & 0.052 & 0.004 & 0.000 & 0.330 & 0.434 \\
blind undercov. & 83 & 0.047 & 0.053 & 0.006 & 0.000 & 0.357 & 0.430 \\
blind undercov. & 131 & 0.049 & 0.054 & 0.004 & 0.000 & 0.315 & 0.434 \\
blind undercov. & 173 & 0.048 & 0.049 & 0.001 & 0.000 & 0.316 & 0.469 \\
blind undercov. & 211 & 0.049 & 0.050 & 0.001 & 0.000 & 0.327 & 0.461 \\
blind undercov. & 257 & 0.051 & 0.055 & 0.004 & 0.000 & 0.309 & 0.406 \\
blind undercov. & 307 & 0.043 & 0.046 & 0.003 & 0.000 & 0.339 & 0.477 \\
blind undercov. & 359 & 0.050 & 0.057 & 0.007 & 0.000 & 0.315 & 0.426 \\
blind undercov. & 401 & 0.047 & 0.053 & 0.005 & 0.000 & 0.334 & 0.398 \\
blind undercov. & 463 & 0.049 & 0.051 & 0.002 & 0.000 & 0.352 & 0.391 \\
blind undercov. & 509 & 0.048 & 0.050 & 0.002 & 0.000 & 0.348 & 0.434 \\
blind undercov. & 571 & 0.051 & 0.052 & 0.001 & 0.000 & 0.362 & 0.418 \\
blind undercov. & 631 & 0.053 & 0.055 & 0.002 & 0.000 & 0.297 & 0.418 \\
blind undercov. & 701 & 0.051 & 0.054 & 0.003 & 0.000 & 0.323 & 0.430 \\
blind undercov. & 761 & 0.051 & 0.057 & 0.006 & 0.000 & 0.316 & 0.438 \\
blind undercov. & 823 & 0.045 & 0.049 & 0.004 & 0.000 & 0.332 & 0.449 \\
blind undercov. & 887 & 0.048 & 0.051 & 0.003 & 0.000 & 0.343 & 0.426 \\
blind undercov. & 947 & 0.048 & 0.054 & 0.006 & 0.000 & 0.322 & 0.449 \\
\addlinespace[2pt]
sparse group04 & 11 & 0.021 & 0.007 & -0.014 & 0.000 & 0.407 & 0.486 \\
sparse group04 & 29 & 0.023 & 0.007 & -0.016 & 0.000 & 0.442 & 0.528 \\
sparse group04 & 47 & 0.020 & 0.007 & -0.013 & 0.000 & 0.413 & 0.472 \\
sparse group04 & 83 & 0.024 & 0.007 & -0.017 & 0.000 & 0.432 & 0.458 \\
sparse group04 & 131 & 0.019 & 0.006 & -0.013 & 0.000 & 0.433 & 0.479 \\
sparse group04 & 173 & 0.024 & 0.006 & -0.018 & 0.000 & 0.399 & 0.472 \\
sparse group04 & 211 & 0.026 & 0.008 & -0.018 & 0.000 & 0.404 & 0.493 \\
sparse group04 & 257 & 0.025 & 0.007 & -0.018 & 0.000 & 0.419 & 0.431 \\
sparse group04 & 307 & 0.025 & 0.007 & -0.018 & 0.000 & 0.412 & 0.403 \\
sparse group04 & 359 & 0.022 & 0.008 & -0.014 & 0.000 & 0.412 & 0.493 \\
sparse group04 & 401 & 0.023 & 0.008 & -0.014 & 0.000 & 0.417 & 0.479 \\
sparse group04 & 463 & 0.020 & 0.005 & -0.014 & 0.000 & 0.435 & 0.507 \\
sparse group04 & 509 & 0.024 & 0.007 & -0.017 & 0.000 & 0.439 & 0.424 \\
sparse group04 & 571 & 0.027 & 0.009 & -0.018 & 0.000 & 0.408 & 0.472 \\
sparse group04 & 631 & 0.022 & 0.007 & -0.015 & 0.000 & 0.399 & 0.451 \\
sparse group04 & 701 & 0.023 & 0.007 & -0.015 & 0.000 & 0.435 & 0.472 \\
sparse group04 & 761 & 0.022 & 0.009 & -0.013 & 0.000 & 0.429 & 0.493 \\
sparse group04 & 823 & 0.026 & 0.005 & -0.021 & 0.000 & 0.410 & 0.500 \\
sparse group04 & 887 & 0.025 & 0.007 & -0.018 & 0.000 & 0.412 & 0.479 \\
sparse group04 & 947 & 0.022 & 0.008 & -0.015 & 0.000 & 0.385 & 0.500 \\
\addlinespace[2pt]
sparse group08 & 11 & 0.023 & 0.004 & -0.019 & 0.000 & 0.409 & 0.497 \\
sparse group08 & 29 & 0.024 & 0.005 & -0.019 & 0.000 & 0.434 & 0.455 \\
sparse group08 & 47 & 0.021 & 0.005 & -0.016 & 0.000 & 0.425 & 0.462 \\
sparse group08 & 83 & 0.022 & 0.004 & -0.018 & 0.000 & 0.417 & 0.455 \\
sparse group08 & 131 & 0.022 & 0.003 & -0.019 & 0.000 & 0.430 & 0.458 \\
sparse group08 & 173 & 0.022 & 0.004 & -0.018 & 0.000 & 0.415 & 0.469 \\
sparse group08 & 211 & 0.024 & 0.004 & -0.020 & 0.000 & 0.405 & 0.455 \\
sparse group08 & 257 & 0.023 & 0.005 & -0.018 & 0.000 & 0.414 & 0.462 \\
sparse group08 & 307 & 0.023 & 0.004 & -0.019 & 0.000 & 0.419 & 0.417 \\
sparse group08 & 359 & 0.022 & 0.005 & -0.017 & 0.000 & 0.419 & 0.472 \\
sparse group08 & 401 & 0.022 & 0.005 & -0.017 & 0.000 & 0.417 & 0.479 \\
sparse group08 & 463 & 0.022 & 0.003 & -0.019 & 0.000 & 0.429 & 0.507 \\
sparse group08 & 509 & 0.023 & 0.004 & -0.019 & 0.000 & 0.428 & 0.476 \\
sparse group08 & 571 & 0.024 & 0.004 & -0.020 & 0.000 & 0.415 & 0.427 \\
sparse group08 & 631 & 0.021 & 0.004 & -0.017 & 0.000 & 0.417 & 0.465 \\
sparse group08 & 701 & 0.021 & 0.004 & -0.017 & 0.000 & 0.434 & 0.490 \\
sparse group08 & 761 & 0.023 & 0.004 & -0.019 & 0.000 & 0.434 & 0.479 \\
sparse group08 & 823 & 0.024 & 0.005 & -0.019 & 0.000 & 0.402 & 0.444 \\
sparse group08 & 887 & 0.025 & 0.004 & -0.021 & 0.000 & 0.419 & 0.476 \\
sparse group08 & 947 & 0.020 & 0.004 & -0.017 & 0.000 & 0.416 & 0.465 \\
\addlinespace[2pt]
\bottomrule
\end{longtable}


\normalsize

\section{Reproducibility Checklist}

\begin{enumerate}[leftmargin=*]
  \item Run the unit tests with \texttt{python3 -m unittest discover -s tests}.
  \item Run the smoke gate with \texttt{./scripts/run\_smoke.sh}.
  \item Regenerate the deep matrix with \texttt{python3 -m experiments.deep\_matrix}.
  \item Regenerate LaTeX inputs with \texttt{python3 scripts/build\_latex\_inputs.py}.
  \item Build the full paper with \texttt{./scripts/build\_latex\_paper.sh}.
  \item Verify that the LaTeX artifact manifest records the PDF hash and page
  count.
  \item Treat \texttt{public/ppo\_grpo\_opd\_long\_horizon.pdf} and the DOCX as
  abridged reports; treat the LaTeX PDF as the full paper artifact.
\end{enumerate}

\section{Next Experiments}

\paragraph{Closed-loop toy training.}
The next toy should train parameterized policies with a shared reward and
compare closed-loop \ppo, \grpo, and corrected \grpo\ variants under equal
generated-token budgets. This will test whether estimator fidelity predicts
policy improvement.

\paragraph{Process rewards.}
Adding process labels or subgoal rewards should narrow the gap between
critic-based and group-relative estimators. This prediction is important
because many practical systems do not rely on terminal reward alone.

\paragraph{Compaction.}
A compaction experiment should split long traces into variable numbers of
sub-traces. The hypothesis is that critic-based individual-rollout objectives
handle this naturally, while group-wise objectives need additional rules for
matching sub-traces across completions.

\paragraph{Distillation hybrids.}
An OPD/OPSD follow-up should use successful RL traces as teacher or privileged
self-teacher material and measure whether dense distillation preserves the
credit assignment gains at lower inference cost.

\section{Replication Protocol for New Benchmarks}

\paragraph{Trace schema.}
A real long-horizon benchmark should record more than a final reward. Each
trajectory should include prompt identity, sampled response identity, token or
step indices, tool calls, tool observations, compaction boundaries, verifier
calls, process labels when available, terminal reward, and any safety or refusal
events. This schema is the minimum needed to compare a group-relative estimator
with a value-model estimator after the fact. Without it, the experiment can
still report final success, but it cannot diagnose where credit was assigned or
lost.

\paragraph{Budget matching.}
The clean comparison is not ``PPO run'' versus ``GRPO run'' in isolation. The
comparison should state generated-token budget, number of prompts, completions
per prompt, verifier calls, critic update steps, policy update steps, and peak
activation or optimizer memory. A critic-based method should be charged for
critic fitting and storage. A group-relative method should be charged for the
extra completions needed to obtain reward contrast. A distillation method should
be charged for teacher inference or self-teacher selection.

\paragraph{Primary measurements.}
For any benchmark with oracle or process annotations, report estimator
correlation with oracle credit, calibrated MSE, sign accuracy, zero-variance
group fraction, and return improvement under equal rollout budgets. For
benchmarks without oracle credit, report held-out process-label agreement,
per-step value calibration where labels exist, and sensitivity to trajectory
length. Final score alone is not enough because a method can improve a short
task while failing the mechanism required for long traces.

\paragraph{Decision rule.}
A public claim that one family is better should require at least three checks:
the mean effect favors the method, the confidence interval excludes practical
near-ties, and an adversarial or counterexample regime is reported. This paper's
toy result follows that rule: it highlights the critic-favorable mean pattern,
the near tie under tiny critic budget, and the group-favorable blind
undercovered case. A larger benchmark should preserve the same discipline.

\paragraph{Audit package.}
The public release should include raw per-seed rows, aggregate tables, plotting
scripts, exact command lines, environment notes, and a short statement of
non-claims. The non-claims matter. For this workstream, the honest claim is that
critic-style temporal information can dominate in observable long traces, while
group-relative normalization remains strong when reward contrast is reliable
and value modeling is weak. That is a research program, not a slogan.

\section{Artifact Provenance and Public Release Notes}

\paragraph{Canonical data source.}
The canonical numeric artifact for this paper is the deep-matrix JSON under
\texttt{results/}. The LaTeX macros, selected result table, axis summary, full
case summary, raw seed rows, and calibrated-error appendix are generated from
that JSON file by the LaTeX input builder. This keeps the manuscript from
silently drifting away from the experiment output.

\paragraph{Rendered artifacts.}
The repository intentionally contains two kinds of public artifacts. The
abridged PDF/DOCX report is designed for quick reading and sharing. The LaTeX
PDF is the full paper artifact, including appendices and raw tables. The two
artifacts should not be compared by page count; they serve different review
contexts.

\paragraph{Manifest interpretation.}
The LaTeX manifest records input hashes, output hashes, the TeX engine, and the
page-count gate. Because generated artifacts are committed with their manifest,
the manifest's recorded git status necessarily describes the generation-time
worktree rather than the final commit that contains the manifest. A clean
checkout can rerun the build script to refresh the manifest and should recover
the same result counts from the canonical JSON.

\paragraph{How to read the claim.}
The public headline should be the cautious one: 16 clear critic-favorable cases,
1 near tie, and 1 clear group-favorable counterexample in this toy estimator
benchmark. The stronger-sounding ``17 mean critic wins'' is true as a
descriptive statistic, but the near-tie is scientifically important. The
counterexample is equally important because it states a boundary condition for
\grpo-style methods.

\bibliographystyle{plainnat}
\bibliography{references}

\end{document}
